\chapter{Basic Testbenches}
\label{ch:tb}

\chapabstract{You can already write a \VHDL{} testbench, and after
\cref{ch:toolkit} you know the \SV{} constructs a verification engineer uses.
This chapter puts the two together, one layer at a time: first a literal
translation, then the race conditions that translation hides, then
self-checking, then file-based vectors, and finally a layered testbench with a
transaction, a generator, a driver, a monitor, a scoreboard and a coverage
collector. At the end you will be able to write a complete class-based
testbench without importing a single line of UVM, and you will understand
exactly which problem UVM was invented to solve.}

Verification is where \SV{} stops being \VHDL{} with different punctuation.
The design half of the language, covered in \refdesignbook, is a
translation exercise. The
verification half is not, because \VHDL{} has no equivalent of most of it:
classes, mailboxes, dynamic threads, constrained randomisation and functional
coverage do not map onto anything you have written before.

\cref{ch:toolkit} taught those constructs one at a time. This chapter is the
other half of the job, and the division of labour between the two is worth
stating once, because it governs everything that follows. \cref{ch:toolkit}
owns the construct: what a mailbox is, what \code{fork ... join\_none} does,
what a covergroup declares. This chapter owns the \emph{use}: why the monitor
puts a transaction into a mailbox rather than calling the scoreboard, why the
driver runs two forked threads that do not know about each other, and why the
covergroup that measures your stimulus tells you nothing about the design. If
a piece of syntax here is unfamiliar, look it up there; if a piece of
structure here seems arbitrary, the answer is in this chapter.

Every testbench, in every language, does the same five things, and the whole
of this chapter is about \emph{organising} those five things better. So we
start by naming them.

% ===========================================================================
\section{The Five Parts of a Testbench}
\label{sec:tb:five}
% ===========================================================================

\index{testbench!structure}A testbench, whether it is fifty lines of \VHDL{}
or fifty thousand lines of \SV{}, consists of exactly five parts:

\begin{enumerate}
\item \textbf{The device under test.} One instance of the thing you are
  verifying, plus whatever wires it needs.
\item \textbf{Clock and reset generation.} Free-running, independent of
  everything else, and usually the first thing you write.
\item \textbf{Stimulus.} Something that puts values on the inputs, in a legal
  order, at legal times.
\item \textbf{Checking.} Something that decides whether the outputs are
  right.
\item \textbf{End of test.} Something that decides when to stop, and reports
  the verdict.
\end{enumerate}

\cref{fig:tb:five} draws them. Nothing in the rest of this chapter adds a sixth
box. The layered testbench of \cref{sec:tb:layered} looks far more complicated
than a \VHDL{} testbench, and it is, but every one of its classes belongs to
box 3 or box 4. The complexity buys separation, reuse and the ability to run
ten thousand random transactions instead of the twelve you thought of by hand.
It does not buy new functionality.

\begin{figure}[htbp]
  \centering
  \begin{tikzpicture}[
      font=\sffamily\scriptsize,
      bx/.style={draw=codeframe,fill=codebg,line width=0.8pt,
                 rounded corners=2pt,minimum width=26mm,minimum height=10mm,
                 align=center,inner sep=3pt},
      dut/.style={bx,fill=svbg,draw=accentalt,line width=1pt,
                  minimum width=26mm,minimum height=16mm},
      ar/.style={-{Stealth[length=2mm]},line width=0.8pt,draw=codekey},
      arg/.style={-{Stealth[length=2mm]},line width=0.8pt,draw=accent,
                  dashed},
      lb/.style={font=\sffamily\tiny\itshape,text=codenumber}
    ]
    \node[bx]  (clk)  at (-4.6, 1.9) {2. clock and reset\\generation};
    \node[bx]  (stim) at (-4.6,-0.3) {3. stimulus};
    \node[dut] (dut)  at ( 0.0, 0.8) {1. DUT};
    \node[bx]  (chk)  at ( 4.6, 0.8) {4. checking};
    \node[bx]  (fin)  at ( 0.0,-2.3) {5. end of test\\and verdict};

    \draw[ar] (clk.east) -- ++(0.9,0) |- (dut.west);
    \draw[ar] (stim.east) -- ++(0.9,0) |- ([yshift=-4mm]dut.west);
    \draw[ar] (dut.east) -- (chk.west);
    \draw[arg] (stim.north east) .. controls (0,2.9) .. (chk.north);
    \node[lb,anchor=south] at (0.2,2.9) {expected values};
    \draw[ar] (chk.south) |- (fin.east);
    \draw[ar] (clk.south) -- (stim.north);
  \end{tikzpicture}
  \caption{The five parts of any testbench. Every construct in this chapter
    belongs to part 3 or part 4.}
  \label{fig:tb:five}
\end{figure}

Keep \cref{fig:tb:five} in mind. When a class-based testbench stops making
sense, the fastest way back is to ask which of the five boxes the class you
are staring at belongs to.

\begin{ruleofthumb}
If you cannot say which of the five parts a piece of testbench code belongs
to, it is in the wrong file.
\end{ruleofthumb}

% ===========================================================================
\section{The Literal Translation}
\label{sec:tb:literal}
% ===========================================================================

Start with something small enough to hold in your head:
\cref{lst:tb:accusv} is an eight-bit synchronous accumulator with an
active-low asynchronous reset and an enable. The \VHDL{} version is the
architecture you would write without thinking about it, and it is not
reproduced here.

\begin{svcode}[caption={An eight-bit accumulator, the DUT for this section},label={lst:tb:accusv}]
`timescale 1ns/1ps

module accu (
  input  logic       clk,
  input  logic       rst_n,
  input  logic       en,
  input  logic [7:0] d,
  output logic [7:0] q
);

  logic [7:0] acc;

  always_ff @(posedge clk or negedge rst_n) begin
    if (!rst_n)   acc <= '0;
    else if (en)  acc <= acc + d;
  end

  assign q = acc;

endmodule : accu
\end{svcode}

Now the testbench. \cref{lst:tb:tbvhdl} is a perfectly ordinary \VHDL{}
testbench for the \VHDL{} accumulator: reset, four accumulate cycles with the
values 10, 20, 30 and 40, then one check that the result is 100.

\begin{vhdlcode}[caption={A conventional \VHDL{} testbench},label={lst:tb:tbvhdl}]
library ieee;
use ieee.std_logic_1164.all;
use ieee.numeric_std.all;

entity tb_accu is
end entity tb_accu;

architecture behavioral of tb_accu is

  constant TCLK : time := 10 ns;

  signal clk   : std_logic := '0';
  signal rst_n : std_logic := '0';
  signal en    : std_logic := '0';
  signal d     : unsigned(7 downto 0) := (others => '0');
  signal q     : unsigned(7 downto 0);
  signal done  : boolean := false;

begin

  dut : entity work.accu
    port map (clk => clk, rst_n => rst_n, en => en,
              d => d, q => q);

  clk <= not clk after TCLK / 2 when not done else '0';

  stim : process is
  begin
    rst_n <= '0';
    wait for 3 * TCLK;
    wait until rising_edge(clk);
    rst_n <= '1';

    for i in 1 to 4 loop
      wait until rising_edge(clk);
      wait for 1 ns;
      en <= '1';
      d  <= to_unsigned(i * 10, 8);
    end loop;

    wait until rising_edge(clk);
    wait for 1 ns;
    en <= '0';

    wait until rising_edge(clk);
    wait for 1 ns;
    assert q = to_unsigned(100, 8)
      report "accumulator wrong: expected 100, got "
             & integer'image(to_integer(q))
      severity error;

    report "test finished" severity note;
    done <= true;
    wait;
  end process stim;

end architecture behavioral;
\end{vhdlcode}

\cref{lst:tb:tbsv} is the line-by-line \SV{} translation. Read it against the
\VHDL{} before reading the commentary.

\begin{svcode}[caption={The same testbench translated literally into \SV{}},label={lst:tb:tbsv}]
`timescale 1ns/1ps

module tb_accu;

  localparam time TCLK = 10ns;

  logic       clk   = 1'b0;
  logic       rst_n = 1'b0;
  logic       en    = 1'b0;
  logic [7:0] d     = '0;
  logic [7:0] q;

  accu dut (.clk(clk), .rst_n(rst_n), .en(en),
            .d(d), .q(q));

  always #(TCLK/2) clk = ~clk;

  initial begin
    rst_n = 1'b0;
    #(3*TCLK);
    @(posedge clk);
    rst_n = 1'b1;

    for (int i = 1; i <= 4; i++) begin
      @(posedge clk);
      #1ns;
      en = 1'b1;
      d  = 8'(i * 10);
    end

    @(posedge clk);
    #1ns;
    en = 1'b0;

    @(posedge clk);
    #1ns;
    assert (q == 8'd100)
      else $error("accumulator wrong: expected 100, got %0d", q);

    $display("[%0t] test finished", $time);
    $finish;
  end

endmodule : tb_accu
\end{svcode}

Six things changed. Each of them deserves a paragraph.

\subsection{The timescale}

\index{timescale}\index{timeunit}\refdesignbook explains
\code{`timescale}, the scoped \code{timeunit}~/ \code{timeprecision}
alternative, and the rule by which a file with no directive silently inherits
whichever one the compiler read last. Read that gotcha again before you write a
testbench, because a testbench is where it does its damage: the clock generator
is the one piece of code whose behaviour depends entirely on what \code{\#10}
means, and a half-period that has quietly become picoseconds looks exactly like
a simulation that hangs. The second failure mode is precision. With
\code{1ns/1ns} the delay \code{\#0.5} is rounded to \code{\#1} without a
warning in most tools, so a \SI{100}{\mega\hertz} clock silently becomes
\SI{50}{\mega\hertz} and every timing number in your report is wrong by a
factor of two.

\begin{ruleofthumb}
\code{`timescale 1ns/1ps} at the top of every file. Deviate only when you are
modelling something sub-picosecond, which in digital verification you are not.
\end{ruleofthumb}

\subsection{\texttt{initial} versus a process with a final \texttt{wait}}

\index{initial block}An \code{initial} block is a \VHDL{} process with no
sensitivity list that ends in \code{wait;} (\refdesignbook). The difference
that matters to a testbench writer is the default. A \VHDL{} process loops back
to the top unless you stop it, so the one-shot stimulus idiom needs an explicit
final \code{wait}; an \code{initial} block runs once and its thread terminates,
so the \emph{forever} case is the one you must ask for, with \code{always} or
with \code{forever} inside an \code{initial}.

\begin{rosetta}
\begin{rleft}
stim : process is
begin
  -- do something once
  a <= '1';
  wait for 10 ns;
  a <= '0';
  wait;              -- never resumes
end process stim;

tick : process is
begin
  -- runs forever
  clk <= not clk;
  wait for 5 ns;
end process tick;
\end{rleft}
\begin{rright}
initial begin
  // do something once
  a = 1'b1;
  #10ns;
  a = 1'b0;
end                  // thread ends here

always begin
  // runs forever
  clk = ~clk;
  #5ns;
end
\end{rright}
\end{rosetta}

Several \code{initial} blocks in the same module all start at time zero and run
concurrently, exactly like several \VHDL{} processes. The order in which they
start is not defined, which is the first hint of the trouble in
\cref{sec:tb:races}.

\subsection{Clock generation}

The \VHDL{} idiom \code{clk <= not clk after 5 ns;} is a concurrent signal
assignment: it re-executes every time \code{clk} changes, so it oscillates on
its own. The \SV{} idiom is procedural.

\begin{rosetta}
\begin{rleft}
-- concurrent assignment
clk <= not clk after TCLK/2;

-- or, in a process
process is
begin
  clk <= '0';
  wait for TCLK/2;
  clk <= '1';
  wait for TCLK/2;
end process;
\end{rleft}
\begin{rright}
// procedural, one statement
always #(TCLK/2) clk = ~clk;

// or, spelled out
always begin
  clk = 1'b0;
  #(TCLK/2);
  clk = 1'b1;
  #(TCLK/2);
end
\end{rright}
\end{rosetta}

Two details. First, \code{clk} must have an initial value: \code{~1'bx} is
\code{1'bx}, so an uninitialised clock stays \code{x} forever and nothing ever
happens. The same trap exists in \VHDL{}, where \code{not 'U'} is \code{'U'},
which is why you already write \code{signal clk : std\_logic := '0';}. Do the
same in \SV{}: \code{logic clk = 1'b0;}. A variable initialiser in a module is
executed once at time zero, before any \code{initial} block.

Second, use a blocking assignment for the clock. This is one of the very few
places where \code{=} is correct in a testbench, and the reason is explained in
\cref{sec:tb:races}: the clock must change in the Active region so that
everything sensitive to the edge sees a clean, single event.

\begin{gotcha}{\texttt{always \#5 clk <= \textasciitilde clk;} is wrong}
With a non-blocking assignment the clock edge happens in the NBA region, the
same region where the design's flip-flops update, and every sampling guarantee
you thought you had is gone. Some tools warn, some do not. Use
\code{clk = \textasciitilde clk}, always.
\end{gotcha}

\subsection{Reset sequencing}

There is nothing new here except syntax, but note the shape of the idiom in
\cref{lst:tb:tbsv}: hold reset asserted for a few clock periods, then release
it \emph{synchronously}, on a clock edge, so that the release has a defined
relationship to the first active edge the design sees. An asynchronous release
at an arbitrary time simulates fine, but every run then starts with a slightly
different alignment, and in a gate-level simulation with real timing checks the
recovery/removal check on the reset pin will fire.

\begin{tipbox}{A reset task, not a reset paragraph}
Put the reset sequence in a task from the first day, even in a flat testbench:
\code{task automatic apply\_reset(int unsigned cycles = 5);} that drives
\code{rst\_n} low, waits \code{cycles} clock edges and releases it
non-blockingly on an edge. When the testbench grows into the layered form of
\cref{sec:tb:layered}, that task becomes \code{cordic\_drv::reset()} unchanged.
Note the default argument value, which \VHDL{} also has, and which is the
cheapest way to keep a task configurable without a second overload.
\end{tipbox}

\subsection{Ending the simulation}

\index{\$finish}A \VHDL{} testbench normally ends by \emph{starvation}: every
process has suspended on a \code{wait} that will never resume, no signal has a
pending transaction, so the simulator has no next event and stops. That is why
\cref{lst:tb:tbvhdl} has the \code{done} signal: without it the clock generator
would keep scheduling events forever.

\SV{} does not rely on starvation. \code{\$finish} stops the simulator
immediately and returns control to the shell; \code{\$stop} suspends it and
drops you into the interactive prompt, which is what you want under a debugger
and never what you want in a regression.

\begin{rosetta}
\begin{rleft}
-- end by starvation:
-- stop the clock, then
-- suspend forever
done <= true;
wait;

-- there is no
-- "stop the simulator now"
-- statement in VHDL-2008
\end{rleft}
\begin{rright}
// end explicitly
$finish;      // exit simulator
$stop;        // break to prompt

// $finish takes a verbosity
// code: 0, 1 (default) or 2
$finish(2);
\end{rright}
\end{rosetta}

\begin{gotcha}{\texttt{\$finish} kills everything, immediately}
\code{\$finish} is not \code{return}: it terminates the whole simulation where
it executes. A monitor halfway through reconstructing a transaction is
discarded, an unflushed scoreboard queue is never checked, and an open file may
lose its last lines. Never call \code{\$finish} from inside a checking thread.
Call it from one place, after an explicit drain, and put anything that must
happen at the very end in a \code{final} block, which is the one thing
\code{\$finish} does run.
\end{gotcha}

\subsection{Printing and asserting}

\index{\$display}\refdesignbook gives the full mapping from \VHDL{}'s
\code{report ... severity} to \code{\$display}, \code{\$info},
\code{\$warning}, \code{\$error} and \code{\$fatal}, and the C-style format
specifiers that go with them. Two consequences of that mapping shape every
testbench in this chapter.

The first is that \SV{} splits printing from counting. \code{\$display} only
prints; the severity tasks print \emph{and} increment the simulator's error
counters, which is what a batch regression reads to decide pass or fail. A
mismatch reported with \code{\$display} is invisible to the regression script.
Report failures with \code{\$error} or \code{\$fatal}, always, and reserve
\code{\$display} for progress and configuration.

The second is \code{\$sformatf}, which formats into a \code{string} instead of
printing. Every \code{convert2string()} method in this chapter is one call to
it, and that is what lets a transaction describe itself in a failure message
written by someone else.

\begin{notebox}{The shape of an immediate assertion}
An immediate assertion has the shape of an \code{if}: the statement after the
condition is the \emph{pass} action, the one after \code{else} is the
\emph{fail} action, and both are optional. Writing a pass action is almost
always a mistake, because it executes on every successful check and floods the
log. Always give a label and a message naming the values involved:
\code{a\_result\_ok : assert (q == 8'd100) else \$error("got \%0d", q);}. The
label appears in the failure message, in the coverage database, and, for
concurrent assertions, in the command that switches the assertion off.
\end{notebox}

At this point you have a working \SV{} testbench that does the same job as the
\VHDL{} one. It is also, as it stands, quietly broken.

% ===========================================================================
\section{Why That Testbench Is Dangerous}
\label{sec:tb:races}
% ===========================================================================

\index{race condition}Look again at the stimulus loop in
\cref{lst:tb:tbsv}. It waits for a clock edge, waits \SI{1}{\nano\second}, then
drives. The \SI{1}{\nano\second} is there because you have written it in every
\VHDL{} testbench you have ever produced. Remove it and see what happens.

\begin{svcode}[caption={A racy stimulus loop. Do not write this},label={lst:tb:racy}]
initial begin
  rst_n = 1'b1;
  for (int i = 1; i <= 4; i++) begin
    @(posedge clk);
    en = 1'b1;          // blocking, on the edge
    d  = 8'(i * 10);    // blocking, on the edge
  end
end
\end{svcode}

Two processes are now sensitive to the same event. The DUT's
\code{always\_ff @(posedge clk)} wants to read \code{en} and \code{d}; the
testbench's \code{initial} block wants to write them. Both are triggered by the
same edge, and both run in the \textbf{Active} region of the same time slot.

\refdesignbook explained what the \LRM{} says about that: processes
scheduled in the Active region execute in an order that is explicitly left
undefined. There is no rule. The simulator may run the testbench first, in
which case the DUT sees the \emph{new} value of \code{en} and accumulates on
this very edge; or it may run the DUT first, in which case the DUT sees the
\emph{old} value and accumulates one cycle later. Both are conforming
implementations.

The observable result is a testbench that produces 100 on one simulator and 60
on another, or 100 today and 60 after an unrelated module changes the compile
order and with it the tool's internal process ordering. There is no warning. In
\VHDL{} this class of bug does not exist, because the language separates the
update phase from the evaluate phase and a signal assignment made during
evaluation is never visible until the next delta cycle. \VHDL{} bought
determinism with the signal/variable distinction; Verilog did not.

\subsection{Remedy (a): drive on the other edge}

The first fix is the one your \VHDL{} instincts reach for: move the driving
away from the edge the design samples.

\begin{svcode}
initial begin
  for (int i = 1; i <= 4; i++) begin
    @(negedge clk);       // half a period away from the DUT
    en = 1'b1;
    d  = 8'(i * 10);
  end
end
\end{svcode}

This works. It is deterministic, it is simple, and for a small testbench with
one clock it is defensible. It has two defects.

The first is that it does not model reality. No real device drives a
synchronous bus on the falling edge. Run this testbench against a gate-level
netlist with back-annotated timing and the stimulus arrives half a period
early, so every setup check passes trivially and the simulation tells you
nothing about whether the interface closes timing. Add a second clock domain,
or a negative-edge-triggered sub-block, and the convention breaks silently.

The second is that it is a convention, not a declaration. It lives in your
head. A colleague adding a new stimulus block six months later has no way of
knowing the rule except by reading every existing block and inferring it.

\subsection{Remedy (b): drive with non-blocking assignments}

The second fix uses the scheduler rather than the waveform.

\begin{svcode}
initial begin
  for (int i = 1; i <= 4; i++) begin
    @(posedge clk);
    en <= 1'b1;           // non-blocking
    d  <= 8'(i * 10);
  end
end
\end{svcode}

Now the assignment is queued in the Active region and applied in the
\textbf{NBA} region. The DUT's \code{always\_ff}, whichever order it runs in,
evaluates the right-hand side of its own non-blocking assignments in Active,
where \code{en} and \code{d} still hold their old values. The design therefore
always sees the old value, on every simulator, deterministically.

This is a genuine fix and it costs nothing. It is the same reason \RTL{} code
uses \code{<=} inside \code{always\_ff}: non-blocking assignment makes the
order in which parallel processes execute irrelevant. Use it everywhere in a
testbench that drives \DUT{} inputs from a clocked context.

\begin{gotcha}{You cannot read back what you just non-blocking-assigned}
\begin{svcode}
en <= 1'b1;
$display("en is %b", en);   // prints the OLD value
\end{svcode}
The assignment has been \emph{queued}, not performed: \code{\$display} runs in
Active, the update happens later in NBA. You already have this reflex for
\VHDL{} signals, but the trap is sharper here, because the same variable could
equally well have been assigned with \code{=}, in which case the print shows
the new value. Use \code{\$strobe}, which prints in the Postponed region after
everything has settled, or read the value after the next clock edge.
\end{gotcha}

\begin{gotcha}{Mixing \texttt{=} and \texttt{<=} on the same testbench signal}
A signal driven with \code{=} in one block and \code{<=} in another is a race
even if neither block is clocked: the blocking write lands in Active, the
non-blocking write lands in NBA of the same time slot, and NBA always wins
regardless of source order. The symptom is a signal that ignores the blocking
assignment entirely and takes the non-blocking value, even when the blocking
line is the later of the two in the source and ran last. Pick one operator per signal:
clock with \code{=}, everything reaching the \DUT{} with \code{<=} or a
clocking block, local bookkeeping with \code{=}.
\end{gotcha}

\subsection{Remedy (c): a clocking block}

Non-blocking assignment removes the race but still leaves timing to
convention: the stimulus changes exactly at the edge, which is a zero
propagation delay that no hardware has. A \emph{clocking block} declares the
timing instead, and it is the right answer.
\index{clocking block}

\refdesignbook introduced the syntax. Here is the real one, from
\file{cordic\_if.sv}, retyped:

\begin{svcode}[caption={The clocking blocks from \file{cordic\_if.sv}},label={lst:tb:cb}]
interface cordic_if (input logic clk);

  import cordic_pkg::*;

  logic   rst_n, req_valid, req_ready, rsp_valid, rsp_ready;
  data_t  req_x, req_y, rsp_x, rsp_y;
  angle_t req_theta;

  clocking drv_cb @(posedge clk);
    default input #1step output #1ns;
    output req_valid, req_x, req_y, req_theta;
    output rsp_ready;
    input  req_ready, rsp_valid, rsp_x, rsp_y;
  endclocking

  clocking mon_cb @(posedge clk);
    default input #1step;
    input req_valid, req_ready, req_x, req_y, req_theta;
    input rsp_valid, rsp_ready, rsp_x, rsp_y;
  endclocking

  modport drv (clocking drv_cb, output rst_n, input clk);
  modport mon (clocking mon_cb, input rst_n, input clk);

endinterface : cordic_if
\end{svcode}

Take the crucial line apart word by word:

\begin{svcode}
default input #1step output #1ns;
\end{svcode}

\textbf{\code{\#1step}} is a unit of time equal to the smallest precision in
the whole simulation, and it exists for exactly one purpose. A sample taken
\code{\#1step} before the clock edge is a sample taken in the
\textbf{Preponed} region, before any process in that time slot has executed.
It is therefore the value the signal had \emph{immediately before} the edge:
the value a real flip-flop would capture. No design process can have changed
it yet, so there is no race, by construction and not by convention.

\textbf{\code{output \#1ns}} means that an assignment to a clocking-block
output does not take effect at the edge. It is scheduled \SI{1}{\nano\second}
later, in the middle of the cycle, which is what a real driver with a
propagation delay does. This removes the second half of the problem: the
stimulus no longer changes at the instant a flip-flop is sampling.

\cref{fig:tb:skew} shows the resulting timing.

\begin{figure}[htbp]
  \centering
  \begin{tikzpicture}[
      font=\sffamily\scriptsize,
      wv/.style={line width=1pt,draw=codekey},
      wd/.style={line width=1pt,draw=accentalt},
      ed/.style={draw=codenumber,dashed,line width=0.6pt},
      an/.style={font=\sffamily\tiny,align=center},
      ar/.style={-{Stealth[length=1.6mm]},line width=0.7pt}
    ]
    % --- clock -------------------------------------------------------------
    \draw[wv] (0,2.6) -- (1,2.6) -- (1,3.2) -- (3,3.2) -- (3,2.6)
              -- (5,2.6) -- (5,3.2) -- (7,3.2) -- (7,2.6) -- (8,2.6);
    \node[anchor=east,font=\sffamily\scriptsize] at (-0.15,2.9) {\texttt{clk}};

    % --- DUT output changes exactly at the edge ---------------------------
    \draw[wv] (0,1.3) -- (1,1.3) -- (1,1.9) -- (5,1.9) -- (5,1.3) -- (8,1.3);
    \node[anchor=east,font=\sffamily\scriptsize] at (-0.15,1.6)
         {\texttt{req\_ready}};
    \node[an,text=codenumber,anchor=north] at (2.9,1.15)
         {driven by the DUT\\at the edge};

    % --- driver output, 1 ns after the edge -------------------------------
    \draw[wd] (0,0.0) -- (1.6,0.0) -- (1.6,0.6) -- (5.6,0.6)
              -- (5.6,0.0) -- (8,0.0);
    \node[anchor=east,font=\sffamily\scriptsize] at (-0.15,0.3)
         {\texttt{req\_valid}};

    % --- edges -------------------------------------------------------------
    \draw[ed] (1,-0.7) -- (1,3.4);
    \draw[ed] (5,-0.7) -- (5,3.4);

    % --- annotations -------------------------------------------------------
    \draw[ar,draw=accent] (0.35,3.9) -- (0.95,3.35);
    \node[an,text=accent,anchor=south east] at (0.7,3.9)
         {\texttt{input \#1step}\\sampled here\\(Preponed)};

    \draw[ar,draw=accentalt] (2.5,-0.55) -- (1.65,-0.05);
    \node[an,text=accentalt,anchor=west] at (2.55,-0.5)
         {\texttt{output \#1ns} -- driven \SI{1}{\nano\second} after the edge};

    \node[an,text=codenumber,anchor=south] at (3.0,3.45)
         {clock period};
    \draw[{Stealth[length=1.4mm]}-{Stealth[length=1.4mm]},draw=codenumber,
          line width=0.5pt] (1,3.42) -- (5,3.42);
  \end{tikzpicture}
  \caption{What \code{default input \#1step output \#1ns} buys you. The
    testbench samples strictly before the edge, so it never sees a value the
    \DUT{} produced on that same edge; it drives strictly after the edge, so
    the \DUT{} never samples a value in the instant it is changing.}
  \label{fig:tb:skew}
\end{figure}

Using it is unremarkable, and \refdesignbook has the direction
convention: a clocking block is written from the testbench's point of view, so
\code{req\_ready}, an output of \code{cordic\_rot}, is an \code{input} here.

\begin{svcode}
vif.drv_cb.req_valid <= 1'b1;      // driven 1 ns after the next edge
vif.drv_cb.req_x     <= t.x0;
do @(vif.drv_cb); while (vif.drv_cb.req_ready !== 1'b1);
vif.drv_cb.req_valid <= 1'b0;      // ready was sampled before the edge
\end{svcode}

Those four lines are the whole valid/ready handshake, and
\cref{sec:tb:layered} uses them unchanged as the body of the driver's
\code{send()} task.

\begin{gotcha}{Two drivers on one signal: the clocking block always wins}
A clocking-block output is driven in the Re-NBA region, after Active and after
NBA. If another part of the testbench also assigns the same signal
procedurally, the clocking block overwrites it in every time slot in which both
happen, regardless of source order, and the symptom is a signal that ignores
half your assignments. A signal listed as a clocking-block \code{output} must
be driven \emph{only} through that clocking block. This is why
\code{cordic\_if} keeps \code{rst\_n} out of \code{drv\_cb} and exposes it
directly on the \code{drv} modport.
\end{gotcha}

\begin{gotcha}{\texttt{vif.drv\_cb.req\_ready} is not \texttt{vif.req\_ready}}
The first is a sample taken before the most recent clock edge; the second is
the live signal, right now. They differ by up to a clock cycle, and a testbench
written half one way and half the other appears to see the same event at two
different times. The symptom is a monitor that reports transactions one cycle
later than the driver thinks it sent them, or a \code{wait} that hangs because
it is watching a live signal that changed in a region the waiting thread cannot
observe. Use the clocking-block reference for everything synchronous, and raw
references only for asynchronous signals such as reset.
\end{gotcha}

\begin{notebox}{\texttt{\#1step} is not \texttt{\#1}}
\code{1step} is a keyword, not a number with a unit. It denotes the smallest
time precision used anywhere in the simulation and may appear only as a
clocking-block input skew, where the \LRM{} defines it to mean the Preponed
region of that time slot. \code{\#1} would mean one time \emph{unit},
typically \SI{1}{\nano\second}, which is a thousand times too coarse.
\end{notebox}

\begin{ruleofthumb}
Drive the \DUT{} only through a clocking block. Sample the \DUT{} only through
a clocking block. Then the question "does my testbench race with the design?"
has the answer "no", permanently, and you stop having to think about it.
\end{ruleofthumb}

% ===========================================================================
\section{Making the Testbench Decide}
\label{sec:tb:selfcheck}
% ===========================================================================

\index{self-checking testbench}The testbench in \cref{lst:tb:tbsv} contains one
assertion, checking one value, at one time. That is not a testbench, it is a
waveform generator with a comment. Going from "run it and look at the waves" to
"the testbench decides" is the single largest quality improvement available to
you, and it costs less effort than the waveform inspection it replaces.

A self-checking testbench needs four things: a model of what the correct answer
is, a comparison with a defined tolerance, a record of how many checks passed
and failed, and a guarantee that it terminates with a verdict even when the
\DUT{} misbehaves.

\subsection{The reference model must describe the specification}

\index{reference model}\index{golden model}The reference model, also called the
golden model or the predictor, computes what the output \emph{should} be. There
is exactly one rule about how to write it, and it is not negotiable:

\begin{ruleofthumb}
Write the reference model from the specification, never from the \RTL{}. A
model derived from the implementation reproduces the implementation's bugs and
proves nothing.
\end{ruleofthumb}

This rule is easy to state and constantly violated, because copying the \RTL{}
is so much quicker. Consider the CORDIC. A tempting reference model would
iterate the same micro-rotations the hardware performs:

\begin{svcode}[caption={A worthless reference model: it is the RTL again},label={lst:tb:badmodel}]
// DO NOT DO THIS
function automatic void bad_predict(input cordic_txn t,
                                    output data_t x_ref, output data_t y_ref);
  idata_t  x = idata_t'(t.x0)      <<< GL;
  idata_t  y = idata_t'(t.y0)      <<< GL;
  iangle_t z = iangle_t'(t.theta)  <<< GL;
  for (int i = 0; i < int'(N); i++)
    if (z >= 0) begin
      x = x - (y >>> i);  y = y + (x >>> i);  z = z - ATAN_ROM[i];
    end else begin
      x = x + (y >>> i);  y = y - (x >>> i);  z = z + ATAN_ROM[i];
    end
  x_ref = sat_narrow(x);  y_ref = sat_narrow(y);
endfunction
\end{svcode}

If the \RTL{} uses the wrong \code{ATAN\_ROM} entry, so does this model, and
they agree. If the \RTL{} folds the quadrant incorrectly, so does this model,
and they agree. If \code{ATAN\_ROM} itself is wrong, both are wrong together.
This model detects only typographical differences between two pieces of code
you wrote on the same afternoon in the same frame of mind, and it will pass a
design that computes something which is not a rotation at all.

Compare the real one from \file{tb\_cordic\_pkg.sv}:

\begin{svcode}[caption={The reference model from \file{tb\_cordic\_pkg.sv}: pure specification},label={lst:tb:goodmodel}]
function automatic void predict(input cordic_txn t,
                                output real x_ref, output real y_ref);
  real th = q_to_real(t.theta, QA);
  real xr = q_to_real(t.x0,    QF);
  real yr = q_to_real(t.y0,    QF);
  x_ref = K_GAIN * (xr * $cos(th) - yr * $sin(th));
  y_ref = K_GAIN * (xr * $sin(th) + yr * $cos(th));
endfunction
\end{svcode}

Six lines, and they say nothing about micro-rotations, guard bits, quadrant
folding, saturation or the FSM. They say what the block computes: a rotation by
$\theta$, scaled by the CORDIC processing gain $K$. Reimplement
\code{cordic\_rot} as a lookup table, as a pipelined unrolled array, or as a
multiply by a rotation matrix, and this model still checks it. That is the test
for whether a reference model is any good: could you swap the \DUT{} for a
completely different implementation of the same specification and keep the
model unchanged?

\begin{notebox}{Three legitimate places for a reference model}
The model does not have to be written in \SV{}. In order of increasing
independence: a \code{function} in the testbench package, as here; a C or
C++ routine called through DPI-C (\refsimulationbook), which is what you
do when the model already exists as part of the firmware; and a Python or
MATLAB script that produced a vector file offline, which is the subject of
\refsimulationbook. The third is the most independent, because the person who
wrote the model was solving a mathematical problem and had never seen your
\RTL{}.
\end{notebox}

\subsection{Tolerance: why zero is the wrong answer}

\index{tolerance}For a control block, a checker compares for exact equality and
that is correct. For a fixed-point arithmetic block, exact equality is almost
always the wrong check, and demanding it produces one of two bad outcomes:
either you weaken the reference model until it matches the \RTL{} bit for bit,
which lands you back in \cref{lst:tb:badmodel}, or you spend a week chasing
"failures" that are one least significant bit of rounding.
\file{tb\_cordic\_pkg.sv} declares:

\begin{svcode}
// Comparison tolerance, in LSBs of the output format.  The theoretical
// bound for N = 14 micro-rotations plus the guard-bit truncation is a
// little under 5 LSB; 12 leaves margin without hiding a real bug.
localparam int TOL_LSB = 12;
\end{svcode}

The number 12 is not arbitrary and it is not a knob you turn until the test
passes. It comes from an error budget, and you must be able to write the budget
down:

\begin{itemize}
\item \textbf{Angle quantisation.} The rotation stops after $N=14$
  micro-rotations, so the residual angle is bounded by
  $\arctan(2^{-13}) \approx \SI{1.2e-4}{\radian}$; times the maximum vector
  magnitude, that is on the order of one external LSB.
\item \textbf{Datapath truncation.} Each micro-rotation shifts right and
  truncates. With $GL=2$ guard bits the accumulated truncation over 14
  iterations is bounded by roughly $N \cdot 2^{-GL}$ internal LSBs, a few
  external LSBs.
\item \textbf{Output narrowing.} \code{sat\_narrow} drops the guard bits with
  an arithmetic shift, which rounds towards $-\infty$: up to one more LSB.
\item \textbf{The reference model's own error.} \code{\$cos} and \code{\$sin}
  are double precision, so this term is negligible here, but in a fixed-point
  model it would not be.
\end{itemize}

The total is a little under five LSB. \code{TOL\_LSB} is set to 12, about two
and a half times the bound. That margin absorbs the corner cases the hand
analysis did not cover, without being so loose that a real functional bug slips
through: a bug that inverts a sign, uses the wrong ROM entry or folds the wrong
quadrant produces an error of \emph{thousands} of LSB, not tens. There is a
wide gap between the noise floor and the smallest real bug, and the tolerance
is placed inside that gap.

\begin{gotcha}{Tolerance is not a place to hide a bug}
A tolerance is a claim about the numerical accuracy of the design and it
belongs in the specification. If you have to raise \code{TOL\_LSB} to make a
test pass, you have either found a bug or discovered that the accuracy claim
was wrong, and someone must decide which, in writing. Raising the number until
the regression turns green converts a failing test into a passing test that
checks nothing. The scoreboard guards against this by tracking
\code{worst\_err} and printing it: if the worst observed error creeps from 4 to
11 LSB between two releases the test still passes, but the report says
something changed.
\end{gotcha}

Here is the comparison itself, from the scoreboard:

\begin{svcode}[caption={Tolerance-based comparison, in LSBs of the output format},label={lst:tb:tol}]
real x_ref, y_ref, ex, ey;
predict(t, x_ref, y_ref);

ex = (q_to_real(t.x, QF) - x_ref) * (2.0 ** QF);   // error in LSB
ey = (q_to_real(t.y, QF) - y_ref) * (2.0 ** QF);
if (ex < 0.0) ex = -ex;
if (ey < 0.0) ey = -ey;
if (ex > worst_err) worst_err = ex;
if (ey > worst_err) worst_err = ey;

n_checked++;
if (ex > real'(TOL_LSB) || ey > real'(TOL_LSB)) begin
  n_failed++;
  $error({"MISMATCH %s\n         got  x=%0d y=%0d\n",
          "         want x=%.1f y=%.1f  (err %.2f / %.2f LSB)"},
         t.convert2string(), t.x, t.y,
         x_ref * (2.0 ** QF), y_ref * (2.0 ** QF), ex, ey);
end
\end{svcode}

Note the braces around the two halves of the format string. \SV{} has no
C-style adjacent string literal concatenation: \code{"a" "b"} is a syntax
error. A long format string is split with the concatenation operator,
\code{\{"a", "b"\}}, which is the same operator used on bit vectors.

Note also that the error is converted into LSBs of the output format before being
compared. That is what makes the tolerance a number a specification can state.
Comparing raw real values against \code{1e-4} works too, but the number means
nothing to a reader and has to be rederived whenever \code{QF} changes.

\begin{tipbox}{Always print both values on a mismatch}
A failure message that says \code{assertion failed} costs an hour. A failure
message that prints the stimulus, the observed value, the expected value and
the error in the natural unit costs a minute. \code{convert2string()} exists
for precisely this, and it is the cheapest thing in the whole testbench.
\end{tipbox}

\subsection{Counting, and the \texttt{final} block}

\index{final block}An error counter is a single integer, incremented on every
mismatch, whose value at the end of simulation determines the verdict. Report
it from a \code{final} block, a \SV{} construct with no \VHDL{} counterpart: it
executes once, at the end, whether the simulation ended by \code{\$finish} or
by running out of events.

\begin{svcode}[caption={A verdict that cannot be missed},label={lst:tb:final}]
int unsigned n_checked = 0, n_failed = 0;
localparam int unsigned MIN_TXN = 50;

final begin
  $display(" transactions checked : %0d", n_checked);
  $display(" mismatches           : %0d", n_failed);
  if      (n_failed != 0)        $display(" *** TEST FAILED ***");
  else if (n_checked < MIN_TXN)  $display(
      " *** TEST FAILED: only %0d of %0d transactions ***", n_checked, MIN_TXN);
  else                           $display(" *** TEST PASSED ***");
end
\end{svcode}

\begin{gotcha}{A \texttt{final} block may not consume time}
\code{final} runs in the last time slot, after the simulator has decided to
stop, and it is a function-like context: no \code{\#}, no \code{@}, no
\code{wait}, no time-consuming task calls. Drain queues and wait for
outstanding transactions \emph{before} \code{\$finish}, and leave the
\code{final} block to print.
\end{gotcha}

\subsection{\texttt{\$error} or \texttt{\$fatal}?}

\index{\$fatal}\index{\$error}\code{\$error} prints, increments the error
count, and \emph{continues}. \code{\$fatal} prints and terminates the
simulation immediately, with an exit code taken from its first argument
(\code{\$fatal(1, "...")} is the conventional form; the argument is a
\code{\$finish}-style verbosity code, and a non-zero value is what your
regression script tests).

The choice is not about severity, it is about whether continuing produces
useful information. Use \code{\$error} for a \emph{data} mismatch: one wrong
result does not prevent the next thousand transactions from being checked, and
the pattern of which ones fail is the most valuable debug information you will
get. A run that stops on the first mismatch tells you one thing; a run
reporting "917 of 1000 failed, all with $\theta < 0$" tells you where the bug
is. Use \code{\$fatal} when the testbench itself can no longer function:
randomisation failed, a configuration file is missing, the \DUT{} has hung.
Continuing would then produce thousands of secondary failures that bury the
real one. \file{tb\_cordic\_pkg.sv} follows exactly this split.

Better than choosing once and for all is a switch: read a
\code{+MAX\_ERRORS=n} plusarg (\cref{sec:tb:run}), default it to zero meaning
"never stop", and call \code{\$fatal} when \code{n\_failed} reaches it. Run the
regression with no limit and debug interactively with \code{+MAX\_ERRORS=1}.

\subsection{The watchdog}

\index{watchdog}A \DUT{} that hangs is not a passing test, but that is what a
naive testbench reports: the driver blocks forever waiting for a
\code{req\_ready} that never arrives, the simulation runs out of events (or
runs until the grid job is killed), and no verdict is printed. Every testbench
needs an independent timeout.

\begin{svcode}[caption={An absolute watchdog},label={lst:tb:watchdog}]
localparam time TIMEOUT = 500us;

initial begin
  #TIMEOUT;
  $fatal(1, "WATCHDOG: no completion after %0t", TIMEOUT);
end
\end{svcode}

That is the crude version, and it is better than nothing. The better version is
an \emph{activity} watchdog: it fires not after a fixed simulated time but
after a number of clock cycles with no progress, so it scales automatically
with the size of the test.

\begin{svcode}[caption={An activity watchdog that resets on progress},label={lst:tb:watchdog2}]
localparam int unsigned MAX_IDLE = 5000;
int unsigned idle = 0, last_checked = 0;

initial forever begin
  @(posedge clk);
  if (sb.n_checked != last_checked) begin
    last_checked = sb.n_checked;
    idle         = 0;
  end else begin
    idle++;                       // not `++idle' inside the condition
    if (idle > MAX_IDLE)
      $fatal(1, "WATCHDOG: %0d cycles with no completed transaction",
             MAX_IDLE);
  end
end
\end{svcode}

The increment is deliberately on a line of its own. \code{else if (++idle >
MAX\_IDLE)} is legal \SV{} and Icarus runs it, but Verilator 5.020 aborts with
an internal error on a pre-increment used inside a condition. Side effects
buried in expressions are hard to read anyway, and in a testbench that must run
on more than one tool they are not worth the two saved characters.

\begin{gotcha}{A watchdog is not optional and it is not a debug aid}
The failure mode a watchdog catches is the one that looks most like success. A
hung simulation in a nightly regression produces no error message, no failing
assertion and no exit code until the batch system kills it; the log then ends
mid-sentence and the summary script records "no result" rather than "failed".
Teams do discover months later that a block was never verified because its
testbench had been timing out silently since the third week. Add the watchdog
on day one, before the first check.
\end{gotcha}

\subsection{The test that passed because it checked nothing}
\label{sec:tb:nocheck}

This is the most dangerous failure mode in verification. Suppose the monitor's
clocking block is mis-declared, or the mailbox between the monitor and the
scoreboard is never connected, or a constraint is unsatisfiable so every
transaction is one the \DUT{} never accepts. In each case \code{n\_failed} is
zero at the end of the run, the verdict is \code{TEST PASSED}, and nothing was
verified. \code{n\_failed == 0} is necessary for a pass; it is not sufficient.
The sufficient condition is:

\begin{svcode}
if (n_failed == 0 && n_checked >= MIN_TXN) ... PASSED
\end{svcode}

\file{tb\_cordic\_pkg.sv} implements the weak version of this in
\code{cordic\_sb::report()}:

\begin{svcode}
if (n_failed == 0 && n_checked > 0) $display(" *** TEST PASSED ***");
else                                $display(" *** TEST FAILED ***");
\end{svcode}

and the strong version in \code{cordic\_env::wait\_drain()}, which knows how
many transactions were sent and fails if fewer came back:

\begin{svcode}[caption={Drain with a bound: the guard against a silent pass},label={lst:tb:drain}]
task automatic wait_drain(int unsigned expected, int unsigned timeout_cyc);
  int unsigned waited = 0;
  while (sb.n_checked < expected && waited < timeout_cyc) begin
    @(vif.mon_cb);
    waited++;
  end
  if (sb.n_checked < expected)
    $fatal(1, "timeout: only %0d of %0d transactions completed",
           sb.n_checked, expected);
endtask
\end{svcode}

\begin{ruleofthumb}
Every testbench must be able to fail. Before you trust a new checker, break the
\DUT{} on purpose -- invert a sign, delete a guard bit -- and confirm that the
test goes red. A checker that has never failed has never been tested.
\end{ruleofthumb}

\begin{tipbox}{The five-minute sanity ritual}
When a testbench first reports PASSED, do all five of these before believing
it:
\begin{enumerate}
\item Comment out one term in the \DUT{} and check that it fails.
\item Set \code{TOL\_LSB} to 0 and check that it fails.
\item Check that \code{n\_checked} is the number you expected, not zero and
  not one.
\item Run with a different \code{-sv\_seed} and check that the stimulus
  actually changed.
\item Look at the waveform for exactly one transaction, and confirm the
  handshake looks like the protocol document.
\end{enumerate}
\end{tipbox}
% ===========================================================================
\section{Vectors From a File}
\label{sec:tb:vectors}
% ===========================================================================

Before the layered testbench, one intermediate step that is often enough on
its own: keep the stimulus and the expected results in a file, and let the
testbench read them. It is how you use a model written in another language
without linking anything, and it is completely reproducible --- the file is
the test.

\Refsimulationbook covers the whole technique, including how to generate the
file from a Python model and how to read it from \SV{}, from \VHDL{} and from
C++. Two things belong here, because they are checking bugs rather than file
format bugs.

\begin{svcode}[caption={Reading a vector file the safe way},label={lst:tb:vecread}]
int fd, n_read;
logic [15:0] a, b, expect_q;

fd = $fopen("vectors.txt", "r");
if (fd == 0) $fatal(1, "cannot open vectors.txt");

n_read = 0;
while (!$feof(fd)) begin
  // $fscanf returns the number of items it matched.  Ignoring it is how a
  // truncated file turns into a test that quietly checks nothing.
  if ($fscanf(fd, "%h %h %h\n", a, b, expect_q) == 3) begin
    drive_and_check(a, b, expect_q);
    n_read++;
  end
end
$fclose(fd);

// The vector count is part of the test.
if (n_read != EXPECTED_VECTORS)
  $fatal(1, "read %0d vectors, expected %0d", n_read, EXPECTED_VECTORS);
\end{svcode}

\begin{gotcha}{A short vector file is a silent pass}
\code{\$readmemh} on a file with fewer lines than the array has entries
fills what it finds and leaves the rest at their initial value, with no
error. \code{\$fscanf} on a malformed line returns a count you did not check.
Both produce a run that reports zero failures because it performed zero
comparisons. Count what you read and assert the count --- this is the same
failure as the one in \cref{sec:tb:nocheck}, arriving by a different route.
\end{gotcha}


% ===========================================================================
\section{The Layered Testbench}
\label{sec:tb:layered}
% ===========================================================================

Everything so far has been one \code{initial} block doing everything. That
works until one of four things happens: you want to reuse the pin-level
protocol code in a second test; you want stimulus and checking to run
concurrently rather than in lockstep; you want random stimulus, so the checker
can no longer know the answer in advance; or the \DUT{} pipelines, so that
several transactions are in flight at once and "drive, wait, check" stops being
a possible shape.

The answer to all four is the same: split the testbench into components with
one responsibility each, and let them communicate by passing \emph{transaction
objects} rather than by sharing wires. \cref{fig:tb:layered} shows the result.
Study it before the code, because every class that follows is one box in that
picture.

\begin{figure}[htbp]
  \centering
  \begin{tikzpicture}[
      font=\sffamily\scriptsize,
      comp/.style={draw=codeframe,fill=codebg,line width=0.9pt,
                   rounded corners=2pt,minimum width=23mm,minimum height=8mm,
                   align=center,inner sep=2pt},
      dutb/.style={comp,fill=svbg,draw=accentalt,line width=1.1pt,
                   minimum width=20mm,minimum height=13mm},
      vifb/.style={draw=codekey,fill=notebg,line width=1pt,
                   rounded corners=1.5pt,minimum width=7mm,
                   minimum height=42mm,align=center},
      mbx/.style={draw=noteframe,fill=notebg,line width=0.7pt,
                  minimum width=12mm,minimum height=5mm,align=center,
                  font=\sffamily\tiny},
      ar/.style={-{Stealth[length=1.9mm]},line width=0.8pt,draw=codekey},
      ard/.style={-{Stealth[length=1.9mm]},line width=0.8pt,draw=accent,
                  dashed},
      lb/.style={font=\sffamily\tiny\itshape,text=codenumber}
    ]
    \node[comp,minimum width=50mm] (test) at (2.0, 4.1)
      {\textbf{test} -- how many, which constraints, which directed cases};

    \node[comp] (gen) at (-1.2, 2.3) {generator /\\sequence};
    \node[mbx]  (m1)  at ( 2.0, 2.3) {\texttt{req\_mbx}};
    \node[comp] (drv) at ( 5.2, 2.3) {driver};

    \node[vifb] (vif) at ( 7.9, 0.2) {};
    \node[font=\sffamily\tiny,rotate=90] at (7.9,0.2)
      {\texttt{cordic\_if} + clocking};
    \node[dutb] (dut) at (10.6, 0.2) {\textbf{DUT}\\\texttt{cordic\_rot}};

    \node[comp] (mon) at ( 5.2,-1.3) {monitor\\(passive)};
    \node[mbx]  (m2)  at ( 2.0,-1.3) {\texttt{obs\_mbx}};
    \node[comp] (sb)  at (-1.2,-1.3) {scoreboard\\+ reference model};
    \node[mbx]  (m3)  at ( 2.0,-2.9) {\texttt{cov\_mbx}};
    \node[comp] (cov) at (-1.2,-2.9) {coverage\\collector};

    \draw[ar] (gen) -- (m1);
    \draw[ar] (m1)  -- (drv);
    \draw[ar] (drv.east) -- (drv -| vif.west);
    \draw[ar] (mon -| vif.west) -- (mon.east);
    \draw[ar] (mon) -- (m2);
    \draw[ar] (m2)  -- (sb);
    \draw[ar] (mon.south) |- (m3.east);
    \draw[ar] (m3)  -- (cov);
    \draw[{Stealth[length=1.9mm]}-{Stealth[length=1.9mm]},line width=0.9pt,
          draw=accentalt] (vif.east) -- (dut.west);

    \draw[ard] (test.west) -| (gen.north);
    \node[lb,anchor=south east] at (-0.2,3.0) {starts, configures};

    \node[draw=accent,dashed,line width=0.9pt,rounded corners=3pt,
          fit=(gen)(drv)(mon)(sb)(cov)(m1)(m3),inner sep=4pt] (env) {};
    \node[font=\sffamily\tiny\bfseries,text=accent,anchor=north west]
      at ([xshift=1mm,yshift=-0.5mm]env.south west) {environment};
  \end{tikzpicture}
  \caption{The layered testbench. Solid arrows carry transaction objects
    through mailboxes; the double arrow is the only pin-level connection in the
    whole picture. \cref{ch:uvm} replaces each box with a UVM base class and
    each mailbox with an analysis port, and changes nothing else.}
  \label{fig:tb:layered}
\end{figure}

Notice what the picture says. Only two components touch a wire: the driver and
the monitor. Everything else works with objects, so everything else is
independent of the pin-level protocol and survives a change to it unchanged.
That is the entire point of the layering.

\subsection{The transaction}

\index{transaction}A transaction is a class holding what goes into the \DUT{}
and what came out of it, for one unit of work. It is a \emph{value}, not a
component: created, filled in, handed on, and eventually garbage collected.

Why a class and not a set of signals? It travels through a mailbox as a single
item, so producer and consumer never have to agree on timing. It carries its
own randomisation constraints, so the rules about what constitutes a legal
request live with the definition of a request. And it can print itself, which
turns a failure message from a puzzle into a diagnosis.

\begin{svcode}[caption={The transaction from \file{tb\_cordic\_pkg.sv}},label={lst:tb:txn}]
class cordic_txn;

  // ---- stimulus ------------------------------------------------------
  rand data_t  x0;
  rand data_t  y0;
  rand angle_t theta;

  // ---- observed result, filled in by the monitor ----------------------
  data_t x;
  data_t y;

  // ---- bookkeeping ----------------------------------------------------
  int unsigned id;
  static int unsigned next_id = 0;

  constraint c_theta {
    theta inside {[-PI_EXT : PI_EXT]};
  }
  constraint c_magnitude {
    x0 inside {[-6000 : 6000]};
    y0 inside {[-6000 : 6000]};
  }
  constraint c_corners {
    theta dist {
      0                 :/ 1,
      PI_EXT            :/ 1,
      -PI_EXT           :/ 1,
      PI_EXT / 2        :/ 1,
      -PI_EXT / 2       :/ 1,
      [-PI_EXT:PI_EXT]  :/ 20
    };
  }

  function new();
    id = next_id++;
  endfunction

  function string convert2string();
    return $sformatf(
      "txn#%0d  x0=%0d (%.4f)  y0=%0d (%.4f)  theta=%0d (%.4f rad)",
      id, x0, q_to_real(x0, QF), y0, q_to_real(y0, QF),
      theta, q_to_real(theta, QA));
  endfunction

endclass : cordic_txn
\end{svcode}

The \code{id} field and the \code{static next\_id} counter matter more than
they look. \code{next\_id} is a \kw{static} member (\cref{sec:toolkit:class}), so
\code{next\_id++} in the constructor gives every transaction a unique serial
number. When transaction 4172 fails in a run of ten thousand, that number is
how you find it in the waveform, in the log and in the coverage database.

The three constraints deserve a note each. \code{c\_theta} states a legality
rule taken straight from the specification: angles outside $\pm\pi$ are not
accepted, and \code{cordic\_rot} has an immediate assertion that fires if you
send one. \code{c\_magnitude} is a \emph{test} decision rather than a legality
rule: it keeps $|(x,y)| \cdot K$ inside the output range so that saturation is
exercised only by directed tests, where the expected behaviour is unambiguous.
\code{c\_corners} uses \code{dist} to bias the distribution, because a uniform
draw over $[-\pi, \pi]$ hits exactly zero, exactly $\pi$ and exactly $\pi/2$
with probability approximately zero, and those are the three angles where the
quadrant-folding logic changes behaviour.

\begin{gotcha}{Reusing one transaction object for a whole run}
This is what the handle-is-a-reference rule of \cref{sec:toolkit:class} looks like at
the place it actually bites you. \code{mbx.put(t)} puts the \emph{handle} into
the mailbox, not a copy of the object. If the generator then does
\code{t.theta = ...; mbx.put(t);} again, both entries point at the same object
and the consumer sees the second value twice. The fix is one line,
\code{cordic\_txn t = new();} \emph{inside} the loop rather than above it;
forgetting it produces a testbench that appears to send $N$ transactions and
actually sends the same one $N$ times, at full speed, reporting a clean pass.
The same rule applies downstream: a scoreboard that stores the handle it was
given, rather than a copy, holds a value that the monitor can still change.
\end{gotcha}

\begin{notebox}{\texttt{convert2string} is not a UVM invention}
The name comes from UVM, where \code{uvm\_object} declares
\code{convert2string()} as a virtual method you override. There is nothing
magical about it: it is an ordinary function returning a \code{string}. Write
it now, and when you move to UVM in \cref{ch:uvm} the method is already there
and the base class simply starts calling it for you.
\end{notebox}

\subsection{The generator}

\index{generator}The generator produces transactions and puts them into a
mailbox. That is all it does. It does not know how a transaction reaches the
pins, how long one takes, or whether the \DUT{} is ready.

\begin{svcode}[caption={Random and directed stimulus generation},label={lst:tb:gen}]
task automatic send_random(int unsigned n);
  repeat (n) begin
    cordic_txn t = new();
    if (!t.randomize()) $fatal(1, "randomization failed");
    req_mbx.put(t);
  end
endtask

task automatic send_directed(data_t x0, data_t y0, angle_t theta);
  cordic_txn t = new();
  t.x0 = x0; t.y0 = y0; t.theta = theta;
  req_mbx.put(t);
endtask
\end{svcode}

\begin{gotcha}{An unchecked \texttt{randomize()} turns a random test directed}
\cref{sec:toolkit:random} explains why the return value must be tested. Here is what
ignoring it costs. An over-constrained transaction keeps whatever values it
had before the failed solve, which for a freshly constructed object is all
zeros. The generator then pushes a thousand identical zero transactions, the
driver drives them, the scoreboard checks them, and the run reports a thousand
passing checks against one vector. Every number in the report looks healthy.
\code{void'(t.randomize());} is worse than the bare call, because it documents
that you meant to ignore the failure. When a solve does fail, use the tool's
constraint debugger rather than deleting constraints at random.
\end{gotcha}

Randomisation changes the shape of the testbench because it decouples
\emph{what} you test from \emph{how many} tests you run. A directed \VHDL{}
testbench contains as many checks as you had patience for. A
constrained-random testbench contains one check and a number that you raise
from 100 to 100\,000 when the regression machine is idle overnight.
\cref{sec:tb:closure} explains how you find out whether the extra runs buy
anything.

The constraint language itself is \cref{sec:toolkit:random}'s subject. What is new
here is where the constraints live: on the transaction class, not in the
generator. That is deliberate. A rule about what constitutes a legal request
belongs with the definition of a request, so that every test inherits it and no
test can forget it, while a rule about what \emph{this} test wants goes in an
inline \code{with \{...\}} clause at the call site. Put a test's preference
into a class constraint and you have made it everybody's preference.

\subsection{The driver}

\index{driver}The driver takes transactions out of the mailbox and turns them
into pin wiggles. It knows the protocol: the handshake, the setup order, the
reset sequence, and nothing else. It does not know what the values mean, does
not check anything, and never looks at the response.

\index{virtual interface}One declaration in \cref{lst:tb:drv} needs a word
before you read the rest of it. \code{virtual cordic\_if vif} is a
\emph{handle} to an interface instance, not an instance: a class object lives
on the heap and has no position in the module hierarchy, so it cannot name
\code{tb\_cordic\_top.intf} the way one module names another. The handle is
\kw{null} until the top-level module hands the real instance to the
constructor, which is the whole of the argument \code{new(intf)} in
\cref{lst:tb:top}. Until \cref{sec:tb:mech} takes the mechanism apart, read
\code{vif.drv\_cb.req\_valid} as ``the \code{req\_valid} signal of whichever
interface instance this driver was given, driven through the \code{drv\_cb}
clocking block''.

\begin{svcode}[caption={The driver from \file{tb\_cordic\_pkg.sv}},label={lst:tb:drv}]
class cordic_drv;

  virtual cordic_if vif;
  mailbox #(cordic_txn) req_mbx;      // from the generator
  int unsigned stall_pct = 25;        // random back-pressure

  function new(virtual cordic_if vif, mailbox #(cordic_txn) req_mbx);
    this.vif     = vif;
    this.req_mbx = req_mbx;
  endfunction

  task automatic reset();
    vif.rst_n              <= 1'b0;
    vif.drv_cb.req_valid   <= 1'b0;
    vif.drv_cb.req_x       <= '0;
    vif.drv_cb.req_y       <= '0;
    vif.drv_cb.req_theta   <= '0;
    vif.drv_cb.rsp_ready   <= 1'b0;
    repeat (5) @(vif.drv_cb);
    vif.rst_n <= 1'b1;
    @(vif.drv_cb);
  endtask

  task automatic send(cordic_txn t);
    vif.drv_cb.req_valid <= 1'b1;
    vif.drv_cb.req_x     <= t.x0;
    vif.drv_cb.req_y     <= t.y0;
    vif.drv_cb.req_theta <= t.theta;
    do @(vif.drv_cb); while (vif.drv_cb.req_ready !== 1'b1);
    vif.drv_cb.req_valid <= 1'b0;
  endtask

endclass : cordic_drv
\end{svcode}

Read \code{send()} with the protocol in mind. The payload and
\code{req\_valid} are driven together through the clocking block, so they
appear \SI{1}{\nano\second} after a clock edge. The loop then advances one
clock at a time and inspects the \emph{sampled} \code{req\_ready}. Because the
sample is taken in Preponed, the value tested at edge $n$ is the value
\code{req\_ready} held during cycle $n-1$, so the loop exits on the first edge
that follows a cycle in which both valid and ready were high. That is the
valid/ready contract exactly. Note also what \code{send()} does \emph{not} do:
it never changes the payload while \code{req\_valid} is high, which is required
by the concurrent assertion \code{a\_req\_stable} inside \code{cordic\_rot}.

The second thread is back-pressure. A response channel that is always ready is
never stalled, so the \DUT{}'s \code{ST\_DONE} state lasts exactly one cycle in
every transaction and any bug in holding the result is invisible. The driver
therefore runs a second, independent process that randomly de-asserts
\code{rsp\_ready}:

\begin{svcode}[caption={Two independent driver threads},label={lst:tb:drvrun}]
task automatic run();
  fork
    forever begin
      cordic_txn t;
      req_mbx.get(t);
      send(t);
    end
    forever begin
      if ($urandom_range(99) < stall_pct) begin
        vif.drv_cb.rsp_ready <= 1'b0;
        repeat ($urandom_range(3) + 1) @(vif.drv_cb);
      end
      vif.drv_cb.rsp_ready <= 1'b1;
      @(vif.drv_cb);
    end
  join_none
endtask
\end{svcode}

\code{fork ... join\_none} (\cref{sec:toolkit:fork}) starts both branches and returns
immediately, so \code{run()} does not block its caller. That is the property
being used here, and it is the reason every component in this testbench
exposes a \code{run()} that starts threads instead of a task that does work:
the environment can start five components in five statements and still have
control afterwards. The two branches are genuinely concurrent and independent,
neither knowing the other exists. There is no way to express this in \VHDL{},
where the number of processes is fixed at elaboration and a subprogram cannot
create one.

\begin{ruleofthumb}
Random back-pressure on every ready signal, from the first test. A protocol bug
that only appears under stall is the most common bug in a valid/ready design
and the least likely to be found by hand-written stimulus.
\end{ruleofthumb}

\begin{tipbox}{Make the stall percentage a knob}
\code{stall\_pct} is a plain public variable, so the test can set
\code{env.drv.stall\_pct = 0} for a throughput measurement, \code{= 90} for a
stress run, and leave it at the default for regressions. Anything a test might
want to vary should be a variable on the component, not a literal buried in a
task.
\end{tipbox}

\subsection{The monitor}

\index{monitor}The monitor watches the pins and reconstructs transactions. It
is \emph{strictly passive}: it never drives anything, and its clocking block
\code{mon\_cb} lists every signal as an input precisely so that it cannot.

The temptation, always, is to have the driver tell the monitor what it sent.
That saves code and destroys the testbench: the checking is then against what
the testbench \emph{intended} to send, not against what appeared on the wires,
so a driver bug becomes invisible, because the same wrong information reaches
the checker from both sides. Valid asserted for two cycles, a payload changed
mid-handshake, a transaction dropped when \code{req\_ready} deasserts: none of
those is detectable by a checker that was told the answer. There are three
further, practical reasons: a passive monitor can
be attached to a gate-level netlist where no driver exists, to the same block
instantiated inside a larger system and driven by real neighbouring logic, and
to an emulation or formal flow, all unchanged.

\begin{svcode}[caption={The monitor from \file{tb\_cordic\_pkg.sv}},label={lst:tb:mon}]
class cordic_mon;

  virtual cordic_if vif;
  mailbox #(cordic_txn) obs_mbx;      // to the scoreboard
  mailbox #(cordic_txn) cov_mbx;      // to the coverage collector

  // Requests accepted but whose result has not come back.
  protected cordic_txn pending[$];

  task automatic run();
    forever begin
      @(vif.mon_cb);
      if (vif.rst_n !== 1'b1) continue;

      // Request accepted -> remember what went in.
      if (vif.mon_cb.req_valid === 1'b1 &&
          vif.mon_cb.req_ready === 1'b1) begin
        cordic_txn t = new();
        t.x0    = vif.mon_cb.req_x;
        t.y0    = vif.mon_cb.req_y;
        t.theta = vif.mon_cb.req_theta;
        pending.push_back(t);
      end

      // Response accepted -> complete the oldest outstanding request.
      if (vif.mon_cb.rsp_valid === 1'b1 &&
          vif.mon_cb.rsp_ready === 1'b1) begin
        cordic_txn t;
        if (pending.size() == 0) begin
          $error("monitor: response with no outstanding request");
        end else begin
          t   = pending.pop_front();
          t.x = vif.mon_cb.rsp_x;
          t.y = vif.mon_cb.rsp_y;
          obs_mbx.put(t);
          cov_mbx.put(t);
        end
      end
    end
  endtask

endclass : cordic_mon
\end{svcode}

The \code{pending} queue is the part worth studying. \code{cordic\_rot} takes
$N+2$ cycles per transaction, and with back-pressure the response can be held
arbitrarily long, so the request going in and the response coming out are
separated in time. The monitor therefore cannot form a complete transaction in
one place: it records the request when the request handshake completes, holds
it in a queue, and completes it when the response handshake completes. For a
deeply pipelined \DUT{} with several transactions in flight this is not an
optimisation, it is the only way the monitor can work at all.

\code{cordic\_txn pending[\$]} is a queue (\cref{sec:toolkit:dyn}), and the choice
of container is a statement about the protocol. A queue with
\code{push\_back}/\code{pop\_front} models \emph{in-order} completion, which
is what \code{cordic\_rot} guarantees. A design that can reorder its responses
needs an associative array indexed by the transaction tag,
\code{cordic\_txn pending[int];}, and nothing else in the monitor changes. If
you pick the wrong one, the monitor will happily pair every response with the
wrong request and the scoreboard will report a design that computes nonsense.

\begin{gotcha}{\texttt{===} and \texttt{!==}, not \texttt{==} and \texttt{!=}}
With \code{==}, a comparison involving \code{x} or \code{z} returns
\code{1'bx}, and \code{if (1'bx)} takes the false branch. Where that bites is
the negative test: \code{if (sig != 1'b1)} is \emph{false} when \code{sig} is
\code{x}, so an unknown signal passes a check meant to catch it. \code{===} and
\code{!==} compare \code{x} as itself and always give a clean 0 or 1. In
\VHDL{} the trap does not exist, because \code{'X' = '1'} is ordinary equality
on an enumeration and is simply \code{false}. Use \code{===} and \code{!==}
everywhere in a monitor.
\end{gotcha}

\begin{gotcha}{A monitor that skips reset must skip it correctly}
\code{if (vif.rst\_n !== 1'b1) continue;} discards everything seen while reset
is asserted. That is right, but it also discards the cycle in which reset
releases, and it does \emph{not} clear \code{pending}. If the \DUT{} can be
reset mid-transaction, the queue still holds the orphaned request afterwards
and it is matched against the first response of the next test, producing a
mismatch that looks like an arithmetic bug and is a testbench bug. Clear the
queue on the reset edge.
\end{gotcha}

\subsection{The scoreboard}

\index{scoreboard}The scoreboard receives observed transactions, computes what
they should have been, compares, and keeps the statistics. Its reference model
and its tolerance were the subject of \cref{sec:tb:selfcheck}; what remains is
the shape of the component.

\begin{svcode}[caption={The scoreboard from \file{tb\_cordic\_pkg.sv}},label={lst:tb:sb}]
class cordic_sb;

  mailbox #(cordic_txn) obs_mbx;
  int unsigned n_checked = 0;
  int unsigned n_failed  = 0;
  real         worst_err = 0.0;

  function new(mailbox #(cordic_txn) obs_mbx);
    this.obs_mbx = obs_mbx;
  endfunction

  task automatic run();
    forever begin
      cordic_txn t;
      real x_ref, y_ref, ex, ey;
      obs_mbx.get(t);
      predict(t, x_ref, y_ref);
      // ... tolerance comparison, exactly as shown earlier ...
    end
  endtask

  function void report();
    $display(" scoreboard: %0d transactions checked, %0d failed",
             n_checked, n_failed);
    $display(" worst error: %.2f LSB (tolerance %0d LSB)",
             worst_err, TOL_LSB);
    if (n_failed == 0 && n_checked > 0) $display(" *** TEST PASSED ***");
    else                                $display(" *** TEST FAILED ***");
  endfunction

endclass : cordic_sb
\end{svcode}

Two structural observations. First, \code{run()} is an infinite loop that
blocks on \code{obs\_mbx.get(t)}. It never terminates and it never polls; the
mailbox suspends the thread until something arrives. This is the standard shape
for every consumer component in a layered testbench. Second, \code{report()} is
a \code{function}, not a task, so it consumes no time and can be called from a
\code{final} block.

For a \DUT{} whose outputs are exact rather than approximate, a scoreboard is
usually a queue of expected transactions, a queue of observed ones, and a
comparison of the two fronts. Here the expected value is instead computed on
demand from the observed stimulus, which works because the model is pure and
immediate. A reference model with state -- a memory model, a protocol with
modes -- must be driven by the same transaction stream in the same order, and
the scoreboard is then also the place where that ordering is enforced.

\subsection{The coverage collector}

\index{functional coverage}\index{covergroup}\cref{sec:toolkit:cov} explains what a
covergroup declares and why functional coverage answers a different question
from code coverage. The component here answers a third question: where in a
testbench does the sampling happen. The scoreboard reports "was the answer
right?"; the coverage collector reports "what did I actually try?", and a
regression can be green for weeks while exercising a third of the design.

\begin{svcode}[caption={The coverage collector from \file{tb\_cordic\_pkg.sv}},label={lst:tb:cov}]
class cordic_cov;

  mailbox #(cordic_txn) cov_mbx;

  angle_t sampled_theta;
  data_t  sampled_x0;

  covergroup cg_stimulus;
    option.per_instance = 1;

    cp_quadrant : coverpoint sampled_theta {
      bins q3_neg = {[-PI_EXT     : -PI_EXT/2 - 1]};
      bins q4_neg = {[-PI_EXT/2   : -1]};
      bins zero   = {0};
      bins q1_pos = {[1           :  PI_EXT/2 - 1]};
      bins q2_pos = {[PI_EXT/2    :  PI_EXT]};
    }
    cp_sign_x : coverpoint sampled_x0 {
      bins neg  = {[-32768 : -1]};
      bins zero = {0};
      bins pos  = {[1 : 32767]};
    }
    x_quadrant_sign : cross cp_quadrant, cp_sign_x;
  endgroup

  function new(mailbox #(cordic_txn) cov_mbx);
    this.cov_mbx = cov_mbx;
    cg_stimulus  = new();
  endfunction

  task automatic run();
    forever begin
      cordic_txn t;
      cov_mbx.get(t);
      sampled_theta = t.theta;
      sampled_x0    = t.x0;
      cg_stimulus.sample();
    end
  endtask

endclass : cordic_cov
\end{svcode}

\begin{gotcha}{A covergroup inside a class is built in the constructor, or not at all}
\cref{sec:toolkit:cov} notes that a covergroup must be constructed. Inside a class
there is a further restriction that catches people: the instance may be created
\emph{only} in that class's \code{function new()}. Forget
\code{cg\_stimulus = new();} and every \code{sample()} call does nothing, the
coverage report shows 0\%, and you spend an afternoon convinced the stimulus is
broken.
\end{gotcha}

Now the distinction the heading promised. \code{cg\_stimulus} covers
\emph{stimulus}: it records what the testbench sent, which tells you whether
the constraints produce the variety you assumed. It says nothing about what the
\DUT{} did internally. A design can receive every quadrant and still never have
its saturation logic activate, never stall in \code{ST\_DONE}, and never see
two requests back to back.

Covering \DUT{} state means sampling signals inside the design, which needs a
second covergroup in a module, usually clocked and usually bound to the \DUT{}
(\cref{sec:tb:sva}):

\begin{svcode}[caption={Covering DUT state rather than stimulus},label={lst:tb:covstate}]
covergroup cg_dut_state @(posedge clk_i);
  cp_state : coverpoint dut.state_q {
    bins idle   = {dut.ST_IDLE};
    bins rotate = {dut.ST_ROTATE};
    bins done   = {dut.ST_DONE};
    bins stalled_in_done = (dut.ST_DONE => dut.ST_DONE);
  }
  cp_sat : coverpoint (dut.x_q[IW-1:GL] > MAX_DATA ||
                       dut.x_q[IW-1:GL] < MIN_DATA) {
    bins saturating = {1'b1};
    bins normal     = {1'b0};
  }
  x_state_sat : cross cp_state, cp_sat;
endgroup
\end{svcode}

The \code{(ST\_DONE => ST\_DONE)} bin is a \emph{transition bin}: it is hit
only when the FSM stays in \code{ST\_DONE} for two consecutive samples, which
happens only when \code{rsp\_ready} is low. If that bin is empty at the end of
a regression, the back-pressure thread never actually stalled anything and you
have not verified the stall path, no matter how green the scoreboard was.

\begin{ruleofthumb}
Stimulus coverage tells you your constraints work. \DUT{}-state coverage tells
you your verification works. Write both, and never sign off on stimulus
coverage alone.
\end{ruleofthumb}

\subsection{The environment}

\index{environment}The environment owns the components, owns the mailboxes,
connects them and starts them. It contains no stimulus and no protocol
knowledge, which is what makes it reusable across every test.

\begin{svcode}[caption={The environment: build, connect, run},label={lst:tb:env}]
class cordic_env;

  virtual cordic_if vif;

  mailbox #(cordic_txn) req_mbx;
  mailbox #(cordic_txn) obs_mbx;
  mailbox #(cordic_txn) cov_mbx;

  cordic_drv drv;
  cordic_mon mon;
  cordic_sb  sb;
  cordic_cov cov;

  function new(virtual cordic_if vif);
    this.vif = vif;
    // "build"
    req_mbx = new();
    obs_mbx = new();
    cov_mbx = new();
    // "connect": the mailboxes are passed to the constructors
    drv     = new(vif, req_mbx);
    mon     = new(vif, obs_mbx, cov_mbx);
    sb      = new(obs_mbx);
    cov     = new(cov_mbx);
  endfunction

  task automatic run();
    fork
      mon.run();
      sb.run();
      cov.run();
    join_none
    drv.run();
  endtask

  function void report();
    sb.report();
    cov.report();
  endfunction

endclass : cordic_env
\end{svcode}

Three steps happen here, in a strict order. \textbf{Build}: every object is
constructed, nothing is connected and nothing runs. \textbf{Connect}: the
mailboxes are handed to the components that use them, which here is done by
passing them to the constructors, so build and connect are textually merged; in
a larger testbench with optional components you separate them, because you
cannot connect two components until both exist. \textbf{Run}: the consumer
threads are started with \code{join\_none} so they run in the background, and
the driver is started last.

Those three steps, done by hand here, are exactly what UVM automates. UVM's
\code{build\_phase} constructs components top-down, its \code{connect\_phase}
wires analysis ports bottom-up, its \code{run\_phase} starts the threads. The
mechanism is different, being a factory, a configuration database and a phasing
engine, but the structure is the one you are reading. When \cref{ch:uvm} introduces
those three method names, you will already know what belongs in each.

\begin{notebox}{Why the phases have to be separate}
Two components cannot be connected until both have been constructed, and a
component cannot be constructed after another component has already started
running and begun looking for it. A single-phase design therefore only works
when the component graph happens to be a tree that you can build in dependency
order by hand. As soon as there is a cycle -- a scoreboard that needs a handle
to the coverage collector and vice versa -- you need construction to be
complete before connection begins. That, and nothing more mysterious, is why
UVM has phases.
\end{notebox}

\subsection{The test}

\index{test}The test is the top of the pyramid: it decides what varies from
run to run. Number of transactions, which constraints are active, which
directed corner cases to add, how much back-pressure, which seed.

\begin{svcode}[caption={The top-level module and one test},label={lst:tb:top}]
`timescale 1ns/1ps

module tb_cordic_top;

  import cordic_pkg::*;
  import tb_cordic_pkg::*;

  logic clk = 1'b0;
  always #5ns clk = ~clk;

  cordic_if intf (.clk(clk));

  cordic_rot dut (
    .clk_i (clk),                  .rst_ni      (intf.rst_n),
    .req_valid_i (intf.req_valid), .req_ready_o (intf.req_ready),
    .req_x_i (intf.req_x),         .req_y_i     (intf.req_y),
    .req_theta_i (intf.req_theta),
    .rsp_valid_o (intf.rsp_valid), .rsp_ready_i (intf.rsp_ready),
    .rsp_x_o (intf.rsp_x),         .rsp_y_o     (intf.rsp_y)
  );

  cordic_env  env;
  int unsigned n_txn = 200;

  initial begin
    if (!$value$plusargs("NTXN=%d", n_txn)) n_txn = 200;

    env = new(intf);
    env.drv.reset();
    env.run();

    // directed corners first: they are the ones you can reason about
    env.send_directed(X0_UNIT,  0,  angle_t'(0));
    env.send_directed(X0_UNIT,  0,  angle_t'(PI_EXT));
    env.send_directed(X0_UNIT,  0,  angle_t'(-PI_EXT));
    env.send_directed(X0_UNIT,  0,  angle_t'(PI_EXT/2));
    env.send_directed(0,        0,  angle_t'(1234));

    // then the random bulk
    env.send_random(n_txn);

    env.wait_drain(n_txn + 5, 200 * (n_txn + 5));
    env.report();
    $finish;
  end

  // independent watchdog
  initial begin
    #5ms;
    $fatal(1, "WATCHDOG: test did not finish");
  end

endmodule : tb_cordic_top
\end{svcode}

The stimulus is in the test, not in the environment, and that separation is not
stylistic. The environment is the part you reuse: the same driver, monitor,
scoreboard and coverage collector serve every test of this block, and serve it
again when the block is instantiated inside a larger system. The stimulus is
different in every test, by definition. Put a \code{send\_random(1000)} call
inside the environment and you have welded one particular test into the
reusable part; the next test then has to accept it or copy the whole
environment.

Note the directed cases coming first. Random stimulus finds bugs you did not
think of; directed stimulus proves the cases you can reason about by hand.
Start with a handful of vectors whose correct answer you can compute on paper,
so that when the random bulk fails you already know the basic path works.

% ===========================================================================
\section{The Mechanics That Hold It Together}
\label{sec:tb:mech}
% ===========================================================================

\cref{ch:toolkit} gave you the constructs the layered testbench is built from:
handles, mailboxes, semaphores, events, dynamic threads and the lifetime
rules that govern them. This section is about the four decisions the testbench
makes with them, and about the one mechanism \cref{ch:toolkit} could not
cover, because it belongs to neither the class world nor the module world but
to the join between them.

\subsection{The virtual interface}

\index{virtual interface}Start with that join, because it is a genuine
problem. \code{cordic\_if} is an \emph{interface instance}:
elaborated, present in the module hierarchy, with a static name
\code{tb\_cordic\_top.intf}. \code{cordic\_drv} is a \emph{class}: created at
run time, living on the heap, with no position in the hierarchy at all. A class
cannot hold a hierarchical reference to something that might not exist, and
there might be two instances of the interface, or none.

A \code{virtual interface} is the bridge. It is a \emph{handle} to an interface
instance, of type \code{virtual cordic\_if}, and it can be stored in a class
member, passed to a function, and assigned at run time.

\begin{svcode}[caption={How the interface handle reaches the classes},label={lst:tb:vif}]
// in the top-level module: a static instance, elaborated at time 0
cordic_if intf (.clk(clk));

// in the class: a handle, null until someone assigns it
class cordic_drv;
  virtual cordic_if vif;
  function new(virtual cordic_if vif, mailbox #(cordic_txn) req_mbx);
    this.vif = vif;             // the handle is copied in here
  endfunction
endclass

// back in the module: the assignment
initial env = new(intf);        // `intf' converts to `virtual cordic_if'
\end{svcode}

The whole mechanism is that one argument. \code{new(intf)} passes the static
interface instance where a virtual interface handle is expected, and the tool
performs the conversion. The environment stores it and passes it on to the
driver and the monitor. Nothing else is needed for a testbench of this size.

\begin{gotcha}{A null virtual interface fails at run time, not compile time}
\code{virtual cordic\_if vif;} initialises to \code{null}. If you forget to
assign it, the first \code{vif.drv\_cb.req\_valid <= 1'b1;} is a null-handle
dereference: a fatal run-time error with a message like "Illegal reference
through a null virtual interface", pointing at a line that looks perfectly
correct. This is one of the two most common first-day errors in a class-based
testbench (the other is forgetting \code{cg = new()}). If a driver appears to
do nothing at all, check the handle first.
\end{gotcha}

\begin{notebox}{Virtual interfaces, modports and the configuration database}
Two refinements you will meet later. A virtual interface can be typed to a
modport: \code{virtual cordic\_if.drv vif;} restricts the handle to the
driver's view, so an accidental write to a monitor-only signal is a compile
error; tool support is good but not universal, which is why
\file{tb\_cordic\_pkg.sv} uses the plain form. And UVM replaces the constructor argument with a lookup in a configuration
database (\cref{ch:uvm}), so that a deeply nested component can find its
interface without every component above it passing the handle down. For four
components, the constructor argument is clearer.
\end{notebox}

\subsection{Why the components talk through mailboxes}

\index{mailbox}The mailbox itself is \cref{sec:toolkit:fork}'s business. What belongs
here is why there is one at all, because the obvious alternative looks
simpler: the monitor could call \code{sb.check(t)} directly and save an object,
a thread and a blocking call.

The reason is that a direct call is a \emph{dependency}. A monitor that calls
\code{sb.check()} needs a handle to the scoreboard, needs the scoreboard to
exist before the monitor starts, and cannot be attached to a testbench that has
no scoreboard, or two of them, or a scoreboard that wants to think for a while
before answering. A monitor that calls \code{obs\_mbx.put(t)} needs none of
those things. It needs a mailbox handle, which it was given at construction,
and it does not care who is on the other end, whether anyone is, or how long
they take. That is the whole of what makes the monitor reusable: it is the
\emph{connection} that is configurable, not the monitor.

Look back at \cref{lst:tb:env} with that in mind: three mailboxes created and
passed to four constructors are the entire wiring diagram of the testbench.
When \cref{ch:uvm} replaces the mailbox with an analysis port it takes the same
idea one step further, letting the monitor broadcast once to any number of
subscribers without naming one of them.

Bounded or unbounded is a real decision, and it is the one mailbox question
this chapter answers rather than \cref{ch:toolkit}. An unbounded mailbox never
blocks the producer, so \code{send\_random(10000)} in \cref{lst:tb:top} returns
almost immediately with ten thousand objects queued in memory. A bounded one
applies natural back-pressure to the generator, which then produces
transactions at the rate the \DUT{} consumes them. For a testbench of this size
either is fine; for a long random run with large transactions, bound it.

\begin{gotcha}{An unbounded mailbox turns a deadlock into a memory leak}
If the driver thread dies, from a null handle, a fatal error or a
\code{disable} that caught more than you meant, a bounded mailbox blocks the
generator almost immediately and you notice. An unbounded one lets the
generator run to completion, so the test proceeds to \code{wait\_drain}, sits
there until the timeout, and reports a drain failure that gives no hint that
the driver is dead. Bounding the request mailbox to four or eight entries makes
a dead consumer visible within a few transactions.
\end{gotcha}

The other two primitives of \cref{sec:toolkit:fork} have narrower jobs here. A
\code{semaphore} belongs where two stimulus threads want the same physical
resource, a named \code{event} where one thread announces a one-off milestone
such as reset completion. Neither appears in \file{tb\_cordic\_pkg.sv}, because
one driver owns the interface and the reset is sequential with the stimulus. If
you reach for a semaphore in a block-level testbench, check first whether the
real answer is a second mailbox.

\subsection{Threads, and where \texttt{disable} bites}

\index{fork!join\_none}The layered testbench uses two of the fork endings from
\cref{sec:toolkit:fork} and no others. \code{join\_none} is what every component's
\code{run()} uses: start a background thread and give control back to the
caller, which is what lets \code{cordic\_env::run()} start three consumers and
then the driver without ever blocking. \code{join\_any} with \code{disable fork} is the timeout
idiom, and it is the one that needs a warning at the point of use.

\begin{gotcha}{\texttt{disable fork} kills every descendant, not just the last fork}
\code{disable fork} terminates all active descendants of the \emph{current}
thread, which includes threads started by a \code{join\_none} several calls
ago. The classic damage: a task uses the timeout idiom, and the
\code{disable fork} silently kills the monitor that an earlier
\code{env.run()} left running. The scoreboard then receives nothing, the test
finishes with \code{n\_checked == 0}, and the naive verdict logic reports a
pass. Contain the damage by wrapping the timeout in its own
\code{fork ... join}, so the \code{disable} only sees that block's children:

\begin{svcode}
fork begin : timeout_scope
  fork
    wait_for_response();
    #1us;
  join_any
  disable fork;          // only kills children of timeout_scope
end join
\end{svcode}
\end{gotcha}

\subsection{Static versus automatic, at the fork}

\index{automatic lifetime}\index{static lifetime}\cref{sec:toolkit:lifetime} states the
rule: a task declared in a module, interface or package has \emph{static}
lifetime unless you write \kw{automatic}, so one copy of its locals is shared
by every concurrent call. The rule is easy to read and easy to forget. Here is
what it does the first time you fork in a loop, which is where a testbench
meets it.

\begin{svcode}[caption={The bug: all four threads print 4},label={lst:tb:staticbug}]
task print_after(int n);      // `int n' is STATIC here
  #(n * 10ns);
  $display("thread %0d at %0t", n, $time);
endtask

initial begin
  for (int i = 1; i <= 4; i++) begin
    fork
      print_after(i);         // four concurrent calls, ONE copy of n
    join_none
  end
end
\end{svcode}

All four calls share the single static \code{n}. The first sets it to 1 and
suspends at the delay, the second sets it to 2 and suspends, and by the time
any of them wakes up \code{n} is 4. The output is four identical lines, all
saying \code{thread 4}, all at the same time, and the testbench appears to have
started one thread instead of four. There are two fixes and you want both.

\begin{svcode}[caption={The fix: automatic task, automatic loop variable},label={lst:tb:staticfix}]
task automatic print_after(int n);   // each call gets its own n
  #(n * 10ns);
  $display("thread %0d at %0t", n, $time);
endtask

initial begin
  for (int i = 1; i <= 4; i++) begin
    automatic int k = i;             // a fresh k per iteration
    fork
      print_after(k);
    join_none
  end
end
\end{svcode}

\kw{automatic} on the task gives every invocation its own arguments and locals;
\code{automatic int k = i;} inside the loop body gives every \emph{iteration}
its own variable to capture the counter into. Both are needed, because without
the second the forked thread reads \code{i} after the loop has advanced and you
get the same symptom for a different reason.

\begin{gotcha}{Forking in a loop without \texttt{automatic}}
The symptom is always the same and always misleading: $N$ threads that all
behave as if they had received the last value. If you ever see "my four drivers
are all sending transaction 4", this is it, and it is worth recognising on
sight because nothing in the log points at the loop. Note that everything
declared inside a \code{class} is already automatic, which is why the layered
testbench is far less exposed to this than the flat one you started with.
\end{gotcha}

\subsection{Ending a test that contains \texttt{forever} loops}

\index{drain}The monitor, the scoreboard and the coverage collector all run
\code{forever}. None of them will ever finish. So what decides that the test is
over?

The answer, in a hand-built testbench, is a two-part protocol: a
\emph{drain} condition and a \emph{timeout}. The drain condition is a
statement about the testbench's own bookkeeping, "every transaction I sent has
been checked", and the timeout is the guarantee that the wait terminates
even when the condition never becomes true. That is precisely
\code{wait\_drain()} in \cref{lst:tb:drain}: it polls \code{sb.n\_checked}
against the number sent, on every clock, up to a cycle limit, and calls
\code{\$fatal} if the limit is reached.

Once \code{wait\_drain()} returns, the test calls \code{report()} and
\code{\$finish}, which terminates every remaining thread including the three
\code{forever} loops. The \code{final} block, if there is one, still runs.

Be careful with the limit. \code{wait\_drain(n, 200*n)} allows 200 cycles per
transaction, which is generous for a CORDIC that takes 16. Add a test that
stalls \code{rsp\_ready} for a thousand cycles, or raise \code{ITER} to 32, and
the same limit becomes marginal: the test starts failing intermittently with a
drain timeout that reads exactly like a \DUT{} bug. Derive the limit from the
design's real latency, or use the activity watchdog of
\cref{lst:tb:watchdog2}.

\begin{notebox}{This is the problem UVM's objections solve}
Every component in a UVM testbench can \code{raise\_objection} when it starts
something that must not be interrupted and \code{drop\_objection} when it is
done; the phase ends when the count reaches zero. That is a distributed version
of \code{wait\_drain}: instead of one place in the test knowing how many
transactions are outstanding, each component states its own condition and the
framework computes the conjunction. \code{set\_drain\_time()} adds a settling
period, and a global timeout makes a stuck objection fail instead of hang. Same
two-part protocol, bookkeeping moved into the library.
\end{notebox}

% ===========================================================================
\section{Assertions In and Around the Testbench}
\label{sec:tb:sva}
% ===========================================================================

\index{assertion!immediate}\index{assertion!concurrent}\cref{sec:toolkit:sva}
covered the syntax of SystemVerilog Assertions, the sampling semantics and the
operator set. This section answers the question that follows: given a working
testbench and a set of properties, where does each one belong. The answer is
not ``in the \RTL{}'', and it is not ``all in one checker file'' either.

\subsection{Immediate assertions: checking code}

An immediate assertion belongs in checking code, and its most valuable use is
one nobody teaches: checking the testbench.

\begin{svcode}
// inside a scoreboard or a monitor
a_no_orphan : assert (pending.size() > 0)
  else $error("response with no outstanding request");

a_range : assert (t.theta >= -PI_EXT && t.theta <= PI_EXT)
  else $fatal(1, "generator produced an illegal angle: %0d", t.theta);
\end{svcode}

The second one is worth a comment. It checks the \emph{testbench}, not the
\DUT{}: if a constraint is wrong and the generator produces an out-of-range
angle, the \DUT{}'s own assertion will fire and you will spend an hour
suspecting the design. An assertion on your own stimulus, at the point where
the stimulus is produced, turns that hour into a line number.

\subsection{Concurrent assertions: protocol properties}

A concurrent assertion belongs wherever a rule about the \emph{protocol} is
stated, which is usually not in the testbench at all. \code{cordic\_rot}
carries one, because the rule is part of its published contract:

\begin{svcode}[caption={The request-stability property from \file{cordic\_rot.sv}},label={lst:tb:prop}]
property p_req_stable;
  @(posedge clk_i) disable iff (!rst_ni)
    (req_valid_i && !req_ready_o) |=>
      $stable(req_x_i) && $stable(req_y_i) && $stable(req_theta_i);
endproperty

a_req_stable : assert property (p_req_stable)
  else $error("cordic_rot: request payload changed while stalled");
\end{svcode}

Read it as a sentence: on every rising edge of \code{clk\_i}, ignoring anything
that happens while reset is asserted, if \code{req\_valid\_i} is high and
\code{req\_ready\_o} is low in this cycle, then in the \emph{next} cycle
(\code{|=>}, the non-overlapping implication) the three payload signals must be
unchanged. That is the valid/ready stability rule, stated once, checked on
every cycle of every simulation.

Checking the same thing procedurally needs a process, a set of shadow
registers, a condition and the discipline to keep all of it in step with the
protocol:

\begin{rosetta}
\begin{rleft}
-- VHDL-2008: a process, by hand
stab : process (clk) is
  variable px : unsigned(15 downto 0);
begin
  if rising_edge(clk) then
    if req_valid = '1' and
       req_ready = '0' then
      assert req_x = px
        report "payload changed"
        severity error;
    end if;
    px := req_x;
  end if;
end process stab;
\end{rleft}
\begin{rright}
// SVA: one property
property p_stable;
  @(posedge clk)
    disable iff (!rst_n)
    (req_valid && !req_ready)
      |=> $stable(req_x);
endproperty
a_stable : assert property
  (p_stable)
  else $error("payload changed");
\end{rright}
\end{rosetta}

\subsection{\texttt{bind}: checkers without touching the RTL}

\index{bind}Assertions inside the \DUT{} are right when they are part of the design
contract, as \code{a\_req\_stable} is. They are wrong when they belong to
verification: they clutter the source, they need an \code{`ifndef SYNTHESIS}
guard, and a designer will eventually delete them. \refdesignbook
introduced \code{bind}, which attaches a module instance to another module from
the outside without editing the target. This is its main use.

\begin{svcode}[caption={A protocol checker bound to the DUT},label={lst:tb:bind}]
// ---- a standalone checker module, in its own file ---------------------
module cordic_checker
  import cordic_pkg::*;
(
  input logic   clk_i, rst_ni,
  input logic   req_valid_i, req_ready_o,
  input logic   rsp_valid_o, rsp_ready_i,
  input angle_t req_theta_i
);

  // valid must never drop before ready
  a_valid_held : assert property (@(posedge clk_i) disable iff (!rst_ni)
      (req_valid_i && !req_ready_o) |=> req_valid_i)
    else $error("req_valid dropped before req_ready");

  // a response must eventually arrive after an accepted request
  a_eventual_rsp : assert property (@(posedge clk_i) disable iff (!rst_ni)
      (req_valid_i && req_ready_o) |-> ##[1:100] rsp_valid_o)
    else $error("no response within 100 cycles");

  // prove that the interesting corner really happened
  c_theta_pi : cover property (@(posedge clk_i) disable iff (!rst_ni)
      (req_valid_i && req_ready_o && req_theta_i == angle_t'(PI_EXT)));

endmodule : cordic_checker

// ---- attach it, from the testbench file, without editing the DUT ------
bind cordic_rot cordic_checker u_chk (
  .clk_i       (clk_i),
  .rst_ni      (rst_ni),
  .req_valid_i (req_valid_i),
  .req_ready_o (req_ready_o),
  .rsp_valid_o (rsp_valid_o),
  .rsp_ready_i (rsp_ready_i),
  .req_theta_i (req_theta_i)
);
\end{svcode}

The \code{bind} statement names a target module \emph{type}, so every instance
of \code{cordic\_rot} anywhere in the design gets a checker, including instances
buried inside a system-level testbench you have not written yet; you can also
bind to a single instance by naming it. The port connections are written in the
scope of the target, so \code{clk\_i} refers to the target's own \code{clk\_i}.

Because the port expressions are elaborated in the target's scope, a
mismatched width is silently truncated or extended exactly as an ordinary port
connection would be. Write every connection by name and check the widths.

\subsection{\texttt{cover property}: proving a corner happened}

\index{cover property}The \code{c\_theta\_pi} statement in \cref{lst:tb:bind}
is a \code{cover property}, and it exists to answer a question the assertions
cannot: did the corner they guard ever happen? If the count is zero, no
assertion about that corner ever got the chance to fail, and a green run means
nothing. This is the cheapest form of coverage in the language and it is
criminally underused:

\begin{svcode}
// did back-pressure ever actually stall a response?
c_rsp_stalled : cover property (
  @(posedge clk_i) disable iff (!rst_ni)
    rsp_valid_o && !rsp_ready_i);

// did two requests ever arrive back to back?
c_back_to_back : cover property (
  @(posedge clk_i) disable iff (!rst_ni)
    (req_valid_i && req_ready_o) ##1 (req_valid_i && req_ready_o));
\end{svcode}

\begin{ruleofthumb}
For every \code{assert property}, ask whether its antecedent ever becomes true.
If you cannot prove it from the log, write the matching \code{cover property}.
\end{ruleofthumb}

% ===========================================================================
\section{Running It}
\label{sec:tb:run}
% ===========================================================================

\index{QuestaSim}A testbench you cannot run reproducibly is not finished. This
section gives the commands for \tool{QuestaSim}, which is what the department
licenses, and one line each for the two other commercial simulators.
\refsimulationbook covers the open-source route, which has different
constraints: no constrained random, no functional coverage, and a compile-then-run
model rather than an interpreted one.

\subsection{The four commands}

\begin{shcode}
# 1. create and map a working library (once per project)
vlib work
vmap work work

# 2. compile.  Order matters: packages before the things that import them.
vlog -sv -work work \
     +incdir+../rtl \
     ../rtl/cordic_pkg.sv \
     ../rtl/cordic_if.sv \
     ../rtl/cordic_rot.sv \
     ../tb/tb_cordic_pkg.sv \
     ../tb/tb_cordic_top.sv

# 3. optimise, keeping visibility for debug
vopt -o tb_opt +acc=rnb work.tb_cordic_top

# 4. run, batch mode, fixed seed
vsim -c -sv_seed 12345 +NTXN=500 -do run.do work.tb_opt
\end{shcode}

Four things in that sequence are not obvious if your only \tool{QuestaSim}
experience is \code{vcom} and \code{vsim} on \VHDL{}. \code{-sv} tells
\code{vlog} to compile as \SV{}
rather than Verilog-2001; without it, \code{logic}, \code{always\_ff} and every
class is a syntax error. \code{vopt} with \code{+acc=rnb} keeps registers, nets
and blocks visible: \tool{QuestaSim} optimises aggressively by default and
removes signals that no longer need to exist, which makes them invisible in the
waveform and unreachable from a \code{bind} or a hierarchical reference. It
costs simulation speed, so a regression that does not need waveforms should
drop it. \code{-c} runs without the GUI, and \code{-sv\_seed} sets the seed for
the constraint solver and for \code{\$urandom}; \code{-sv\_seed random} picks a
different one each run and prints it, which is what a nightly regression wants.
Finally, \code{+NTXN=500} is a \emph{plusarg}: a command-line argument the
simulation itself can read.

\subsection{Plusargs}

\index{plusarg}\index{\$value\$plusargs}A plusarg is how you change a test's
behaviour without recompiling. \SV{} provides two system functions.

\begin{svcode}[caption={Reading the command line from inside the simulation},label={lst:tb:plusargs}]
int unsigned n_txn      = 200;
int unsigned stall      = 25;
string       vector_file = "";
bit          verbose    = 0;

initial begin
  void'($value$plusargs("NTXN=%d",   n_txn));
  void'($value$plusargs("STALL=%d",  stall));
  void'($value$plusargs("VEC=%s",    vector_file));
  verbose = $test$plusargs("VERBOSE");
  $display("config: NTXN=%0d STALL=%0d VEC=%s VERBOSE=%0b",
           n_txn, stall, vector_file, verbose);
end
\end{svcode}

\code{\$test\$plusargs} returns 1 if the named argument appeared on the command
line at all. \code{\$value\$plusargs} looks for an argument starting
\code{+NTXN=}, parses the rest with the given format, writes it into
\code{n\_txn} and returns 1; if the argument is absent it returns 0 and leaves
the variable alone, which is why initialising the variable at its declaration
is essential.

\begin{gotcha}{\texttt{\$value\$plusargs} leaves the variable untouched on failure}
If the plusarg is missing, or present but unparseable, the target variable
keeps whatever it had. Declare it with a sensible default on the same line, and
print the resulting configuration at the start of every run. A run whose log
does not state its own configuration cannot be reproduced, and a regression
result you cannot reproduce is worth very little.
\end{gotcha}

\begin{notebox}{UVM's plusargs are the same mechanism}
\code{+UVM\_TESTNAME=my\_test}, \code{+UVM\_VERBOSITY=UVM\_HIGH} and
\code{+uvm\_set\_config\_int=...} are ordinary plusargs read with
\code{\$value\$plusargs} inside the UVM library, exactly as in
\cref{lst:tb:plusargs}. There is no simulator support for UVM specifically; it
is a package that reads the same command line you can read.
\end{notebox}

\subsection{A minimal \texttt{.do} file}

\begin{tclcode}
# run.do -- batch execution with a non-zero exit code on failure
onerror {quit -code 3}

# log everything recursively so a failure can be debugged from the WLF
log -r /*

run -all

# TESTSTATUS: 0 = no errors, 1 = warning, 2 = error, 3 = fatal
quit -code [expr {[coverage attribute -name TESTSTATUS -concise] > 1 ? 1 : 0}]
\end{tclcode}

The last line matters more than it looks. A simulation that printed a hundred
\code{\$error} messages still exits with status 0 unless you make it do
otherwise, and a regression script that only checks the exit code records it as
a pass. Either translate the simulator's error count into an exit code as
above, or have your script grep the log for your own \code{*** TEST FAILED ***}
marker. Do both.

\subsection{Waveforms: the dump file is not the viewer}
\label{sec:tb:waves}

\index{waveform}\index{VCD}\index{GTKWave}\index{Surfer}Two separate things are
involved in looking at a waveform, and confusing them is the reason students
ask why their \code{.vcd} will not open. One is the \emph{dump file}, written
by the simulator, which records who changed and when. The other is the
\emph{viewer}, a separate program that draws it.

\tool{QuestaSim} writes its own format, a \file{.wlf} file, and shows it in its
own wave window. \code{log -r /*} in the \code{do} file above is what fills it;
\code{add wave -r /*} is what puts signals on the screen. Run interactively
with \code{vsim -voptargs="+acc=rnb"} and the whole loop stays inside one tool,
which is the fastest way to debug and the reason the department licenses it.

The portable route is a \index{\$dumpfile}\code{\$dumpfile}/\code{\$dumpvars}
pair in the testbench itself. It is plain IEEE 1800, so it works in every
simulator, \tool{QuestaSim} included, and it writes \emph{VCD}: the IEEE 1364
text format that every tool on earth can read. Gate it behind a plusarg,
because a VCD is enormous:

\begin{svcode}
initial if ($test$plusargs("DUMP")) begin
  $dumpfile("tb_cordic.vcd");
  $dumpvars(1, tb_cordic_top.dut);   // this scope only, not everything
end
\end{svcode}

The first argument to \code{\$dumpvars} is a depth: \code{1} means the named
scope only, \code{2} adds its children, and \code{0} means everything below it,
recursively. \code{\$dumpvars(0, tb\_cordic\_top)} on a toy accumulator costs
nothing. The same line on a real design, run for a few million cycles, produces
a file measured in gigabytes, and you will spend longer waiting for the viewer
to index it than you spent on the simulation. Name the scope you care about.

A VCD opens in \tool{GTKWave}, the long-standing open-source viewer, and in
\tool{Surfer}, which is newer, has a modern interface and runs in a browser as
well as natively. \tool{QuestaSim} will also read a VCD if somebody hands you
one. Where the flow allows it, prefer \emph{FST}, a compressed binary format
that is far smaller than VCD and far faster to open; \tool{Verilator} writes it
with \code{-{}-trace-fst} (\refsimulationbook) and both viewers read it. For
anything longer than a few thousand cycles, FST is what you want.

\begin{ruleofthumb}
Dump selectively and dump late. A plusarg-gated \code{\$dumpvars} on one scope,
turned on only for the seed that failed, beats a nightly regression that writes
a gigabyte per run and is never looked at.
\end{ruleofthumb}

\subsection{Coverage, and the other simulators}

\begin{shcode}
# collect coverage: same flags on vlog, vopt and vsim
vlog -sv +cover=bcefsx -work work ../rtl/*.sv ../tb/*.sv
vopt -o tb_opt +cover=bcefsx +acc=rnb work.tb_cordic_top
vsim -c -coverage -sv_seed random work.tb_opt \
     -do "run -all; coverage save run1.ucdb; quit"
vcover merge merged.ucdb run*.ucdb
vcover report -details -html -htmldir covhtml merged.ucdb
\end{shcode}

\code{+cover=bcefsx} asks for branch, condition, expression, FSM, statement and
extended toggle coverage. Those are the \emph{code} coverage metrics of
\cref{sec:toolkit:cov}, computed from the \RTL{} without any help from you. The
functional coverage from \cref{lst:tb:cov} lands in the same UCDB, and the
merge report shows both. Read them as answers to different questions.

The other two commercial simulators use a single compile-and-run command.
\tool{Cadence Xcelium}: \code{xrun -sv -access +rwc -seed 12345 +NTXN=500
*.sv}. \tool{Synopsys VCS}: \code{vcs -sverilog -full64 -debug\_access+all -o
simv *.sv} followed by \code{./simv +ntb\_random\_seed=12345 +NTXN=500}. The
plusarg spelling for the seed is the only genuinely tool-specific thing in this
chapter.

\begin{tipbox}{Put all of this in a Makefile on day one}
Nobody remembers \code{+acc=rnb} at eleven at night. A four-target Makefile
(\code{comp}, \code{sim}, \code{wave}, \code{cov}) that hides the flags is
twenty minutes of work and it is the difference between a testbench that gets
run and one that gets run once. \refsimulationbook gives a complete one.
\end{tipbox}

% ===========================================================================
\section{The Coverage-Closure Loop}
\label{sec:tb:closure}
% ===========================================================================

\index{coverage closure}\cref{ch:toolkit} names this loop in a sentence. It is
worth a section, because it is the daily working rhythm of the whole activity
and because every step of it has a way of going wrong. A constrained-random
testbench is not a test, it is a \emph{machine for producing tests}. Running it
once proves almost nothing. Running it correctly is a loop:

\begin{enumerate}
\item \textbf{Run with $N$ different seeds.} Ten seeds of a thousand
  transactions is a better use of an hour than one seed of ten thousand,
  because a bug that depends on a rare ordering is more likely to appear in ten
  independent sequences.
\item \textbf{Merge the coverage databases} from all the runs into one.
\item \textbf{Look at the holes.} Which bins are empty? Which crosses have zero
  hits? Which \code{cover property} never fired?
\item \textbf{Classify each hole.} Reachable but never produced by the
  constraints: adjust the constraints or add a directed test. Unreachable by
  construction: exclude it with \code{ignore\_bins} and a comment saying why.
  Reachable and \emph{should not be}: you have found a bug in the
  specification or in your understanding of it, which is the most valuable
  outcome available.
\item \textbf{Repeat} until the remaining holes are all excluded with reasons.
\end{enumerate}

What this looks like in practice, honestly, is less tidy. The first run
typically reaches 60 to 80\% functional coverage almost immediately, because
the constraints were written by someone thinking about the interesting cases.
The next 15\% takes a day of adjusting distributions. The last few percent are
individually painful: each is a corner the constraints cannot reach without
being rewritten, and the usual resolution is a small directed test that forces
it, added to the test list with a comment.

Two failure modes are worth naming, because both feel like progress. The first
is chasing coverage instead of bugs: it is entirely possible to reach 100\%
functional coverage on a design with a serious bug, if the covergroups describe
the stimulus you sent rather than the behaviour you needed. Coverage measures
what you tried; the scoreboard measures whether it worked, and a green coverage
report with a weak scoreboard is worse than a red one because it stops people
looking. The second is excluding holes to close the report.
\code{ignore\_bins} exists for genuinely unreachable states, and every use of
it should carry a comment a reviewer can check. An exclusion added the week
before tape-out, with no comment, is how a bug ships.

\begin{ruleofthumb}
Coverage tells you where you have not looked. It never tells you that the
design is correct. Only the checker does that, and only if you have proved the
checker can fail.
\end{ruleofthumb}

\begin{tipbox}{A regression script in five lines}
\begin{shcode}
for seed in $(seq 1 20); do
  vsim -c -sv_seed $seed +NTXN=1000 \
       -do "run -all; coverage save cov_$seed.ucdb; quit" work.tb_opt \
    | tee log_$seed.txt
done
grep -l "TEST FAILED" log_*.txt
vcover merge merged.ucdb cov_*.ucdb && vcover report -details merged.ucdb
\end{shcode}
That is the whole methodology. Everything a commercial regression manager adds
is bookkeeping on top of these three commands.
\end{tipbox}

\begin{notebox}{Where this chapter stops and \cref{ch:uvm} starts}
Nothing in \file{tb\_cordic\_pkg.sv} needs UVM. It is about 470 lines and does
everything described in this chapter. What UVM adds is standardisation: the
same six components with agreed names and base classes, a factory so a test can
substitute one component without editing the environment, a configuration
database so a handle can reach a deeply nested component, phasing, and
objections. Every one of those solves a problem you have now met by hand. That
is why this chapter comes first.
\end{notebox}

% ===========================================================================
\section{Checklist}
\label{sec:tb:checklist}
% ===========================================================================

\textbf{Translating a \VHDL{} testbench}

\begin{itemize}
\item \code{`timescale 1ns/1ps} at the top of every file, or \code{timeunit} in
  every module and package. Never inherit it.
\item A \VHDL{} process ending in \code{wait;} becomes an \code{initial} block.
  A process that loops forever becomes \code{always} or \code{forever}.
\item \code{clk <= not clk after 5 ns;} becomes \code{always \#5ns clk =
  \textasciitilde clk;}. Blocking assignment, and initialise \code{clk} at
  declaration.
\item \code{report ... severity note} becomes \code{\$display}; \code{severity
  error} becomes \code{\$error}; \code{severity failure} becomes
  \code{\$fatal(1, ...)}.
\item \code{assert c report m severity error} becomes \code{a\_name: assert (c)
  else \$error(m)}. Always label it, always give a message with values in it.
\item End with \code{\$finish}, not by starving the event queue.
\end{itemize}

\textbf{Never do this}

\begin{itemize}
\item Never drive a \DUT{} input with \code{=} from a block triggered by the
  same clock edge the \DUT{} samples on. Use \code{<=}, or better, a clocking
  block.
\item Never mix \code{=} and \code{<=} on the same signal.
\item Never call \code{randomize()} without testing its return value.
\item Never let the driver hand transactions to the monitor or the scoreboard.
  The monitor reconstructs from the pins or the testbench proves nothing.
\item Never write a reference model by copying the \RTL{}.
\item Never raise a tolerance to make a test pass without recording why.
\item Never declare a testbench task without \code{automatic}, and never fork
  in a loop without copying the loop variable into an \code{automatic} local.
\item Never trust a \code{TEST PASSED} that does not also state how many
  transactions were checked.
\end{itemize}

\textbf{Always do this}

\begin{itemize}
\item Drive and sample through a clocking block with \code{default input
  \#1step output \#1ns}.
\item Use \code{===} and \code{!==} in a monitor.
\item Give every transaction a unique \code{id} and a
  \code{convert2string()}.
\item Add a watchdog before the first check, and prefer one that times out on
  lack of progress.
\item Assert a minimum transaction count as part of the pass condition.
\item Apply random back-pressure to every ready signal from the first test.
\item Write both stimulus coverage and \DUT{}-state coverage, and a
  \code{cover property} for every assertion whose antecedent might never occur.
\item Bind protocol checkers rather than editing the \RTL{}.
\item Break the design on purpose once, and confirm the testbench goes red.
\item Print the run's full configuration, including the seed, in the first
  lines of the log.
\item Gate \code{\$dumpfile}/\code{\$dumpvars} behind a plusarg and name the
  scope. \code{\$dumpvars(0, top)} on a real design writes gigabytes.
\end{itemize}

\textbf{Structure}

\begin{itemize}
\item Transaction: data plus constraints plus \code{convert2string()}. No
  behaviour.
\item Generator: makes transactions, puts them in a mailbox. Knows nothing
  about pins.
\item Driver: mailbox to pins. Knows the protocol only.
\item Monitor: pins to mailbox. Passive, and independent of the driver.
\item Scoreboard: reference model plus comparison plus statistics.
\item Coverage: what did we actually exercise?
\item Environment: build, connect, run. No stimulus.
\item Test: stimulus, constraints, counts, seed. The only thing that changes
  between runs.
\end{itemize}
